% concatenated from FKG_RI.tex
\documentclass[12pt,a4paper]{article}
% \documentclass[12pt]{article}
\usepackage{amsmath}
\usepackage{amssymb}
\usepackage{accents}

\usepackage{mathtools}
\usepackage{amsthm}
\usepackage[usenames]{color}

\addtolength{\textwidth}{1cm}

\usepackage{graphicx}

\usepackage[numbers,sort&compress]{natbib}
\usepackage{doi}   % prints DOI nicely
\usepackage{url} 

\allowdisplaybreaks

\newtheorem{theo}{Theorem}%[section]
\newtheorem{theorem}[theo]{Theorem}
\newtheorem{lem}[theo]{Lemma}
\newtheorem{lemma}[theo]{Lemma}
\newtheorem{df}[theo]{Definition}
\newtheorem{prop}[theo]{Proposition}
\newtheorem{proposition}[theo]{Proposition}
\newtheorem{cor}[theo]{Corollary}
\newtheorem{corollary}[theo]{Corollary}
\newtheorem{rem}[theo]{Remark}
\newtheorem{remark}[theo]{Remark}
\newtheorem{definition}[theo]{Definition}
\newtheorem{example}[theo]{Example}

%%% to make the labels visible
%  \usepackage{showkeys}

%\newcommand{\qed}{\hfill$\Box$\par\medskip\par\relax}

%%% For indicators. This is used in the following way: 
%%% \1{x=0} produces {\mathbf 1}{\{x=0\}}.
\newcommand{\1}[1]{{\mathbf 1}{\{#1\}}}

\newcommand{\vr}{\varrho}

\newcommand{\muo}{\mu^\omega}
\newcommand{\nuo}{\nu^\omega}

\newcommand{\eps}{\varepsilon}
\newcommand{\Z}{{\mathbb Z}}
\newcommand{\N}{{\mathbb N}}

\newcommand{\C}{{\mathbb C}}
\newcommand{\B}{{\mathsf B}}

\newcommand{\am}{{\mathfrak A}}
\newcommand{\cl}{{\mathfrak C}}

\newcommand{\Sph}{{\mathbb S}}
\newcommand{\V}{{\mathcal V}}
\newcommand{\K}{{\mathcal K}}
\newcommand{\M}{{\mathcal M}}
\newcommand{\LL}{{\mathcal L}}
\newcommand{\G}{{\mathcal G}}
\newcommand{\D}{{\mathcal D}}
\newcommand{\HH}{{\mathcal H}}
\newcommand{\R}{{\mathbb R}}
\newcommand{\RR}{{\mathcal R}_\omega{}}
% \newcommand{\tRR}{{\tilde{\mathcal R}}}
\newcommand{\CC}{{\mathfrak C}}
\newcommand{\cCC}{{\mathcal C}}
\newcommand{\OO}{{\mathcal O}}
\newcommand{\ZZ}{{\mathcal Z}}
\newcommand{\DD}{{\mathcal D}}
\newcommand{\BB}{{\mathcal B}}
\newcommand{\T}{{\mathcal T}}
\newcommand{\NN}{{\mathcal N}}

\newcommand{\Ttr}{{\mathcal T}^{\mathrm{tr}}}
% \newcommand{\Tsp}{{\mathcal T}^{\mathrm{sp}}}
\newcommand{\Tsp}{{\mathcal T}}

\newcommand{\PPP}{{\mathcal P}}
\newcommand{\EE}{{\mathcal E}}
\newcommand{\F}{{\mathcal F}}

\newcommand{\RI}{\mathop{\mathrm{RI}}}
\newcommand{\BRI}{\mathop{\mathrm{BRI}}}

\newcommand{\s}{{\widehat S}}

\newcommand{\gs}{{\mathfrak S}}

\newcommand{\I}{{\mathcal I}}
\newcommand{\J}{{\mathcal J}}
\newcommand{\JJ}{{\mathfrak J}}

\newcommand{\FF}{{\mathfrak F}}

\newcommand{\Q}{{\mathbb Q}}
\let\phi=\varphi
\newcommand{\A}{{\mathcal A}}

\newcommand{\PP}{{\mathbf P}}

\newcommand{\bP}{\mathbf{P}}
\newcommand{\bbP}{\mathbb{P}}
\newcommand{\E}{{\mathbb E}}

\newcommand{\hpsi}{{\widehat\psi}}
%\newcommand{\ttau}{{\tilde T^\vr}}
\newcommand{\hphi}{{\widehat\varphi}}

\newcommand{\bga}{{\bar\gamma}}
\newcommand{\tX}{{\widetilde X}}
\newcommand{\tW}{{\widetilde W}}
\newcommand{\tV}{{\tilde V}}
\newcommand{\tK}{{\tilde K}}
\newcommand{\txi}{{\tilde\xi}}

\newcommand{\tQ}{{\tilde Q}}
\newcommand{\tT}{{\tilde T}}
\newcommand{\diam}{{\mathop{\mathrm{diam}}}}
\newcommand{\8}{{\infty}}
\newcommand{\nn}{{\nonumber}}

\newcommand{\he}{\widehat{e}}
\newcommand{\hL}{\widehat{L}}

%%% equality in law + convergence
\newcommand{\eqlaw}{\stackrel{\text{\tiny law}}{=}}
\newcommand{\eqsetas}{\stackrel{\text{\tiny a.s.}}{=}}
\newcommand{\convlaw}{\stackrel{\text{\tiny law}}{\longrightarrow}}

%\newcommand{\om}{\text{\boldmath ${\omega}$}}

\newcommand{\IQ}{{\mathbb Q}}
\newcommand{\tIQ}{{\widetilde{\mathbb Q}}}
\newcommand{\hIQ}{{\widehat{\mathbb Q}}}
\newcommand{\IP}{{\mathbb P}}
\newcommand{\IE}{{\mathbb E}}
\newcommand{\Cor}{\mathop{\mathrm{Cor}}}

\newcommand{\distTV}{\mathop{dist}\nolimits_{\mathit{TV}}}

\newcommand{\hP}{\widehat{P}}
\newcommand{\tP}{\widetilde{P}}
\newcommand{\hE}{\widehat{E}}
\newcommand{\tE}{\widetilde{E}}

\newcommand{\gep}{\varepsilon} 

\newcommand{\p}{\mathbb{P}}


\newcommand{\bbS}{\mathbb{S}}

\newcommand{\bL}{\mathbb{L}}

\newcommand{\ind}{{\mathbf{1}}} %il vaudrait mieux mathbbm (mais quel package est-ce?)


%%% not-so-wide hat
\DeclareMathSymbol{\widehatsym}{\mathord}{largesymbols}{"62}
\newcommand\lowerwidehatsym{%
  \text{\smash{\raisebox{-1.3ex}{%
    $\widehatsym$}}}}
\newcommand\fixwidehat[1]{%
  \mathchoice
    {\accentset{\displaystyle\lowerwidehatsym}{#1}}
    {\accentset{\textstyle\lowerwidehatsym}{#1}}
    {\accentset{\scriptstyle\lowerwidehatsym}{#1}}
    {\accentset{\scriptscriptstyle\lowerwidehatsym}{#1}}
}
 
\newcommand{\hW}{\widehat{W}}
\newcommand{\hZ}{\widehat{Z}}
\newcommand{\tZ}{\widetilde{Z}}

% \newcommand{\EE}{{\mathfrak E}}
\newcommand{\UU}{{\mathcal U}}
\newcommand{\VV}{{\mathsf V}}

\newcommand{\supp}{\mathop{\mathrm{supp}}}
\newcommand{\essinf}{\mathop{\mathrm{ess\,inf}}}

\newcommand{\near}
{\stackrel{{\scriptstyle \omega}}{\leftrightarrow}}
%\newcommand{\conv}{\mathop{\rm Conv}{\mathcal D}}

\newcommand{\capa}{\mathop{\mathrm{cap}}}
\newcommand{\Cov}{\mathop{\mathrm{Cov}}}
\newcommand{\Var}{\mathop{\mathrm{Var}}}
\newcommand{\hcapa}{\mathop{\widehat{\mathrm{cap}}}}
\newcommand{\hm}{\mathop{\mathrm{hm}}\nolimits}
\newcommand{\hhm}{\mathop{\widehat{\mathrm{hm}}}\nolimits}
\newcommand{\Es}{\mathop{\mathrm{Es}}\nolimits}

\newcommand{\dist}{\mathop{\mathrm{dist}}}

\newcommand{\htau}{\widehat{\tau}}

\newcommand{\tS}{{\widetilde S}}
\newcommand{\tN}{{\widetilde N}}
\newcommand{\tL}{{\widetilde\LL}}
\newcommand{\ttL}{{\widetilde L}}
\newcommand{\tH}{{\widetilde H}}

\newcommand{\ttau}{{\tilde\tau}}

% \newcommand{\hS}{{\hat S}}

\newcommand{\intl}{\int\limits}


%MACROS POUR PARTIE 2 : CHAPITRE GENERAL RI
\newcommand{\wleft}{\overset{\leftarrow}{w}}
\newcommand{\wright}{\overset{\rightarrow}{w}}
\newcommand{\wcenter}{\bar{w}}
\newcommand{\Out}{\mathrm{Out}}
\newcommand{\In}{\mathrm{Ins}}
\newcommand{\Hin}{\mathrm{Hin}}



\newcommand{\fc}{\color{blue}} %%  corrections de FC
\newcommand{\nc}{\color{black}}

\newcommand{\harm}{\textrm{harm}}


\setcounter{tocdepth}{2}%print only sections and subsections in table of contents


\title{Random interlacements on transient weighted graphs: 0-1 laws and FKG inequality}
%{Corrigendum and addendum to: ...}
\author{Orph\'ee Collin$^{1}$}


\begin{document}

\maketitle

{\footnotesize 
\noindent $^{~1}$Technische Universit\"at Wien, Austria
\\
\noindent e-mail:
\texttt{orphee.collin@tuwien.ac.at}
}



\begin{abstract}
We study some properties of the random interlacement model on a transient weighted graph, which was introduced by A. Teixeira in \cite{Teixeira2009}. We give a simple proof of the FKG-property and discuss the occurrence of several 0-1 laws for non-local events. We show in particular a 0-1 law for some increasing non-local events, without any assumption.
\\[.3cm]\textbf{AMS subject classifications (2010 MSC):}
60K35, %Interacting random processes; statistical mechanics type models; percolation theory
82B43, %Percolation
60F20, %Zero-one laws
60J10. %Markov chains (discrete-time Markov processes on discrete state spaces)
\\[.3cm]\textbf{Keywords:} random interlacement model, percolation theory, 0-1 laws
\\[.3cm] This work is licensed under a Creative Commons Attribution | 4.0 International licence (CC-BY 4.0, https://creativecommons.org/licenses/by/4.0/).
\end{abstract}

 

\tableofcontents



\section{Overview}


We consider an irreducible and reversible Markov chain $(X_n)_{n\ge 0}$ on a countable state space. We use the framework of electrical networks: we consider a simple undirected graph $\mathcal{G}=(\mathbb{V}, \mathcal{E})$ and a collection of positive weights (or conductances) $\mathbf{a}=(a_{x,y})_{(x,y)\in E}$ (with $a_{x, y}=a_{y,x}$), and define on $\mathbb{V}$ the Markov chain $X=(X_n)_{n\ge 0}$ by the following transition probabilities: with $a_x:=\sum_{z\sim x} a_{x,z}$,
\begin{equation}
  p_{x,y}:=
  \begin{cases}
  \frac{a_{x,y}}{a_x}, & \text{ if } (x,y)\in \mathcal{E}, \\
  0 &  \text{ else}.
  \end{cases}
\end{equation}
We furthermore assume that $\mathcal{G}$ is connected and locally finite, and that $X$ is transient.

Following the work of A. Teixeira in \cite{Teixeira2009} (generalizing the construction made in \cite{Sznitman2010} for the SRW in $\Z^d$ in transient dimensions), we consider the random interlacement model built upon this transient weighted graph. The model consists of a random collection of biinfinite trajectories of $X$, which we will refer to as the \emph{interlacement process}: it is indeed a (Poisson) point process on the space of transient nearest-neighbor trajectories on $\mathcal{G}$. 
It is in particular interesting to focus on the \emph{interlacement set} $\I$, defined as the union of the ranges of all the trajectories. This random subset $\I$ of $\mathbb{V}$ can be viewed as a dependent percolation set. It is caracterised by the following property: for every finite subset $K$ of $\mathbb{V}$,
\begin{equation}\label{eq:car_prop_I}
  \bbP(\I \cap K = \emptyset)= \exp(- \capa(K)),
\end{equation}
where $\capa(K)$ denotes the capacity of $K$ with respect to $X$.
%The parameter $u$ is an intensity parameter that can be tuned to add more and more trajectories to the interlacement process. For us, it will not pay an important role, and we will most of the time work with $u=1$ (see ??).


In \cite{Teixeira2009}, it was shown that the interlacement set $\I$ satifies the FKG inequality. Here, we point out that this relies in a direct way on the poissonian nature of the interlacement process, and we slightly extend the scope of the FKG inequality.

Next, we study the occurence of a 0-1 law for this model. Contrarily to the case of the SRW on $\Z^d, d\ge 3$, %(or more generally a transitive weighted graph), 
there is no natural candidate for an ergodic set of measure-preserving transformations. Our approach to the 0-1 law rather relies on \emph{non-local} events, that is, events that \emph{do not depend on the configuration inside any finite region} (in a sense that will be defined). 

We establish this 0-1 law if $X$ is tail-trivial - that is, when the $\sigma$-algebra of invariant events is trivial with respect to $X$. %if $X$ satisfies the Liouville property  DANS LA V2
We also give a criterion for the occurence of the 0-1 law when this $\sigma$-algebra is purely atomic.
We further provide a quantitative criterion for the validity of the 0-1 law. 
Additionally, we show that a weaker 0-1 law, concerning events that rely only on one tail of the trajectories, holds in full generality. 
We finally show a 0-1 law for increasing events that are non-local in a slightly stronger sense.

%We conclude by discussing a variety of examples to keep in mind during the proofs, and by presenting some applications of our 0-1 laws. %DANS LA V2


\section{The model}


\subsection{Some notations}


As announced, we consider a simple undirected connected and locally finite graph $\mathcal{G}=(\mathbb{V}, \mathcal{E})$. In the sequel, subsets of $\mathbb{V}$ will be refered to as ``vertex sets''. 
%For any vertex set $K$, let us denote its inner boundary: 
%$$\partial K :=\{x\in K : \exists y \notin K, x\sim y\}.$$
We denote by $\mathcal{P}(\mathbb{V})$ the set of all vertex sets, and by $\mathcal{P}_f(\mathbb{V})$ the set of finite vertex sets. They are endowed with the inclusion order. %=\{K \subset V, K \text{ finite}\}$. 
An exhaustion of $\mathbb{V}$ is an increasing sequence $(K_n)_{n\ge 0}$ of finite subsets of $\mathbb{V}$ such that $\bigcup_{n\ge 0} K_n =\mathbb{V}$. 
For $f: P_f(\mathbb{V})\to \R$ and $l\in \R \cup\{-\infty, \infty\}$, 
we say that $f$ converges towards $l$, 
and write $f(K)\underset{K\to\mathbb{V}}\longrightarrow l$, or $\lim_{K\to \mathbb{V}} f(K) = l $,
if for every exhaustion $(K_n)_{n\ge 0}$ of $\mathbb{V}$, $f(K_n)\underset{n \to\infty}\longrightarrow l$. 
We observe that if $f$ is monotonous, then $\lim_{K\to \mathbb{V}} f(K)$ necessarily exists.


%??? QUESTION: Voir si je définis comme Teixeira, distance de graphe, boule, symbole $\sim$... 

\medskip

Furthermore, we consider a family $\mathbf{a}=(a_{x,y})_{(x,y)\in \mathcal{E}}$ of positive conductances on the graph $\mathcal{G}$ and the induced Markov chain $X=(X_n)_{n\ge 0}$ on $\mathbb{V}$, having the following transition probabilities:
\begin{equation}
  p_{x,y}:=
  \begin{cases}
  \frac{a_{x,y}}{a_x}, & \text{ if } (x,y)\in \mathcal{E}, \\
  0 &  \text{ else}.
  \end{cases}
\end{equation}
with $a_x:=\sum_{z\sim x} a_{x,z}$. 
For every vertex $x$, let $P_x$ denote the law of $X$ started from $x$ and $E_x$ the associated expectation. 
We observe that $(a_x)_{x\in \mathbb{V}}$ is an invariant measure for $X$. Throughout the article, we work under the assumption that $X$ is transient.
For every vertex set $K$, we define the two following stopping times:
\begin{equation}
\begin{aligned}
  \tau_K & = \inf \{n\geq 0 : X_n \in K \} \qquad \text{ (hitting time of } K), \\
  \tau^+_K & = \inf \{n\geq 1 : X_n \in K \}  \qquad\text{ (return time to } K),
\end{aligned}
\end{equation}
as well as the random time:
$$\lambda_K := \sup \{n\ge 0 : X_n \in K\}  \qquad \text{ (time of the last visit in } K),$$
We use the conventions $\inf \emptyset = \infty$ and $\sup \emptyset = -\infty$.
We also define the inner escape boundary of $K$:
\begin{equation}
  \partial_X K :=\{x \in K : P_x[\tau_K^+=\infty]>0\}.
\end{equation}%\subset \partial K.
We observe that, by transience, for any $x\in K$, $P_x$-almost surely, $\lambda_K<\infty$ and $X_{\lambda_K}\in \partial_X K$. 

\medskip

We now introduce the notions of equilibrium measure, capacity and harmonic measure of a finite vertex set $K$. The equilibrium measure is the following measure on $\mathbb{V}$, supported on $\partial_X K$: % (defined by its value on singletons)
\begin{equation}
  e_K(x) := 
  \begin{cases}
    a_x P_x[\tau_K^+=\infty] & \text{if } x\in  K,\\ 
    0 & \text{otherwise}.
  \end{cases}
\end{equation}
The total mass of $e_K$ is called the capacity of $K$ and denoted:
\begin{equation}
  \capa(K):=e_K(\mathbb{V})= \sum_{x\in K} a_x P_x[\tau_K^+=\infty].
\end{equation}
Finally, the harmonic measure of $K$ (or normalized equilibrium measure) is:
\begin{equation}
  \harm_K(\cdot):= \frac{e_K(\cdot)}{\capa(K)}\, .
\end{equation}
Crucial to us is the following property, which derives elementarily from the reversibility of $X$.
%We have the following crucial properties, which justify the term equilibrium. 
%GIVE A REFERENCE (cite location in teixeira : NONE FOUND...).
For every finite vertex sets $K$ and $L$ with $K\subset L$, we have
\begin{equation}\label{eq:consistency_eq_meas}
  \sum_{x\in L} e_L(x) P_x[X_{\tau_K}\in \cdot, \tau_K<\infty] = e_K(\cdot),
\end{equation}
so that, denoting $P_{\harm_L}(\cdot)=\sum_{x\in \mathbb{V}} \harm_L(x) P_x(\cdot)$, we have
\begin{equation}
  P_{\harm_L}[\tau_K<\infty ]= \frac{\capa(K)}{\capa(L)} 
\end{equation}
and
\begin{equation}
  P_{\harm_L}[X_{\tau_K} \in \cdot |\tau_K<\infty] = \harm_K(\cdot) .
\end{equation}


\subsection{Interlacement process}\label{sec:def_interl_process}

We now present the construction of the interlacement process, refering to \cite{Teixeira2009} for details.

We consider the set of transient nearest-neighbor trajectories in $\mathcal{G}$ indexed by $\N \cup \{0\}$ (respectively, by $\Z$):
\begin{equation}
\begin{aligned}
    W^+=\{ w=(w_n)_{n\ge 0} : \quad  & \forall n\ge 0, w_n \sim w_{n+1} \text{ and }  \\
    & \forall x \in \mathbb{V}, \{n \ge 0 :w_n=x\} \text{ is finite}\},
\end{aligned}
\end{equation}
\begin{equation}
\begin{aligned}
    W=\{ w=(w_n)_{n\in \Z} : \quad  & \forall n\in \Z, w_n \sim w_{n+1} \text{ and }  \\
    & \forall x \in \mathbb{V}, \{n \in \Z:w_n=x\} \text{ is finite}\}.
\end{aligned}
\end{equation}
We endow $W^+$ (respectively, $W$) with the $\sigma$-algebra $\mathcal{W^+}$ (respectively, $\mathcal{W}$) generated by the coordinate maps $w \mapsto w_n$, for $n\ge 0$ (respectively, for $n\in \Z$). Note that the laws $P_x, x\in \mathbb{V}$, are probability measures on $(W^+, \mathcal{W}^+)$.


\smallskip

For every finite vertex set $K$, we denote by $W_K$ the set of trajectories in $W$ that hit $K$ and we consider the measure $\nu_K$ on $(W, \mathcal{W})$ caracterized by:
\begin{equation}
\nu_K\left( \left\{ w\in W: (w_{-n})_{n\ge 0} \in A, w_0=x, (w_n)_{n\ge 0}  \right\} \right) 
 = P_x[A|\tau_K^+=\infty] e_K(x) P_x[B], 
\end{equation}
for every $x\in \mathbb{V}$ and $A, B \in \mathcal{W^+}$. This measure has total mass $\capa(K)$ and is supported by the subset of $W_K$ consisting of the trajectories that enter in $K$ at time $0$.


We may then consider a Poisson point process with intensity measure $\nu_K$. For clarity, let us describe how this Poisson point process can be sampled:
\begin{itemize}
\item sample a number $\xi$ as a Poisson variable of parameter $\capa(K)$,
\item independently for each $i=1,  \dots, \xi$, sample a point $x_i$ under $\harm_K$,
\item independently for each $i=1, \dots, \xi$, attach to $x_i$ a trajectory that enters in $K$ through $x$, at time 0, as follows: independently,
\begin{itemize}
  \item for the negative indices, sample a trajectory under $P_x[ \cdot | \tau_K^+=\infty]$ and apply the function $n\mapsto -n$ to the indices,
  \item for the positive indices, sample a trajectory under $P_x[ \cdot ]$.
\end{itemize}
\end{itemize}
Now, let $K\subset L$ be two finite vertex sets. Consider the Poisson point process with intensity measure $\nu_L$. Restrict it to trajectories that hit $K$, and reindex these trajectories so that they enter $K$ at time 0. It follows from \eqref{eq:consistency_eq_meas} that the obtained point process is Poisson, with intensity $\nu_K$; see Theorem 2.1 in \cite{Teixeira2009} for details.

This consistency property allows us to extend the construction to the whole of $\mathcal{G}$: this gives the random interlacement model on $(\mathcal{G}, \mathbf{a})$. Let us construct it formally. To begin with, due to the reindexation in the consistency property, we choose to look at trajectories in $W$ up to time-translation. We define the equivalence relation:
\begin{equation}
  w\sim w' \quad  \text{ if and only if } \quad \exists k\in \Z, (w_{n+k})_{n\in \Z}=(w_n')_{n\in \Z} .
\end{equation}
Denote by $W^*$ the set of equivalence classes of $W$ under this equivalence relation and by $\pi^*:W\to W^*$ the associated projection. Denote by $\mathcal{W}^*$ the $\sigma$-algebra on $W^*$ inherited from $\mathcal{W}$. For each finite vertex set $K$, denote $W_K^*= \pi^*(W_K)$.

Consider now the measure $\nu$ on $(W^*, \mathcal{W}^*)$ caracterised by the fact that for any finite vertex set $K$, the measure $\nu\ind_{W_K^*}$, i.e, the restriction of $\nu$ to trajectories (seen up to time-shift) that hit $K$, is equal to $\nu_K \circ (\pi^*)^{-1}$. 
One way to construct the measure $\nu$ is to fix an exhaustion $(K_n)_{n\ge 0}$ with $K_0=\emptyset$ and set 
$$\nu = 
\left(\sum_{n\ge 1} \nu_{K_n} \ind_{W_{K_n} \setminus W_{K_{n-1}}} \right)
\circ (\pi^*)^{-1}.$$
Then we consider a Poisson point process $(W^*, \mathcal{W}^*)$ with intensity measure $\nu$. Formally, we define the space:
\begin{equation}\label{eq:def_state_Omega}
  \Omega =  \left\{ \omega = \sum_{i\ge 1} \delta_{w_i^*} ,  w_i^* \in W^* : \, \forall K \in \mathcal{P}_f(\mathbb{V}), \, \omega(W_K^*)<\infty \right\}.
\end{equation}
We endow $\Omega$ with the $\sigma$-algebra $\mathcal{A}$ generated by the maps $\omega \mapsto \omega(D)$ for $D\in \mathcal{W}^*$ 
and we consider on $(\Omega, \mathcal{A})$ the law $\bbP$ of a Poisson point process on $(W^*, \mathcal{W}^*)$ with intensity measure $\nu$. We will refer to $\omega =\sum_{i\ge 1} \delta_{w_i^*} $ (under $\bbP$) as the interlacement process, while 
\begin{equation}
  \I:=\bigcup_{i\ge 1} \text{Range}(w_i^*), \qquad  \V = \mathbb{V} \setminus \I 
\end{equation}
(with obvious definition for $\text{Range}(w_i^*)$) are respectively the interlacement set and the vacant set. 
We recover the caracteristic property of $\I$, announced in \eqref{eq:car_prop_I}, by observing that $\I\cap K=\emptyset$ is equivalent to $\omega(W_K^*)=0$ and that $\omega(W_K^*)$ is distributed as a Poisson variable with parameter $\nu(W_K^*)=\nu_K(W_K)=\capa(K)$. 

\begin{remark}\label{rem:full_interlacement_process}
  It is possible to multiply the intensity measure $\nu$ by a positive real $u$; or equivalently to multiply the family $\mathbf{a}$ of conductances by $u$. 
  We thus get a family of interlacement processes indexed by $u$. 
  There is a natural way to construct all these models together through an increasing coupling: 
  consider a Poisson point process $\sum_{i\ge 1} \delta_{(w_i^*, u_i)}$ on $\mathcal{W}^*\times [0, \infty)$ with intensity measure $\nu \otimes d u$; 
  then, for every $u>0$, set the interlacement process at level $u$ to be  $\omega^u:=\sum_{i\ge 1 : u_i\le u} \delta_{w_i^*}$.
  As soon as $u\le v$, $\omega^u$ is dominated by $\omega^v$. 
  In this paper, we will mostly work with $u$ fixed and equal to 1, and therefore disregard this coupling. 
\end{remark}


\section{Theorems}

\subsection{FKG inequality}

% merely a remark


In his seminal paper, A.-S. Sznitman asked whether the interlacement set in $\Z^d$ satifies the FKG inequality (see \cite{Sznitman2010}, Remark 1.6). 
This question was answered positively by A. Teixeira in \cite{Teixeira2009} in the general setup of transient Markov chains, with a detailed proof. 
Here, we point out that this property holds simply because Poisson point processes satisfy in general the FKG inequality. 
We thus show that the interlacement process $\omega$, and not only the interlacement set $\I$, satifies the FKG inequality under $\mathbb{P}$.

\smallskip

Let us recall some definitions.
%\begin{definition}
Let $(\mathcal{X}, \mathcal{F}, \le )$ be a mesurable space equipped with an order. 
A measurable subset $A$ of $\mathcal{X}$ is said to be increasing if 
\begin{equation}
\forall x, y  \in \mathcal{X}, \qquad  ( x \in A \text{ and } x \le y ) \implies y \in A\, .
\end{equation}
If $(\mathcal{X}, \le)$ and $(\mathcal{Y} , \le)$ are two ordered sets, a function $f: \mathcal{X} \to \mathcal{Y}$ is said to be non-decreasing if
\begin{equation}
\forall x, y \in \mathcal{X}, \qquad x \le y \implies f(x) \le f(y)\, . 
\end{equation}
A probability measure $P$ on $(\mathcal{X}, \mathcal{F})$ is said to satisfy the FKG property if for every increasing events $A$ and $B$ we have
\begin{equation}
P[A \cap B ] \ge P[A] \times P[B]\, ,
\end{equation}
or, equivalently, that for every measurable nonnegative and nondecreasing measurable functions $f, g : \mathcal{X} \to \R$, we have
\begin{equation}
E[fg] \ge E[f] \times E[g]\, ,
\end{equation}
where $E$ denotes the expectation under $P$, and $\R$ is equipped with its Borelian $\sigma$-algebra and its usual order. 
An $\mathcal{X}$-valued random variable is said to satisfy the FKG inequality if its law does.
%\end{definition}

\smallskip

The following lemma is an elementary fact.

\begin{lemma}\label{lem:f(X)_FKG}
Let $(\mathcal{X},\mathcal{F} , \le)$ and $(\mathcal{Y}, \mathcal{G} , \le)$ be two mesurable spaces, equiped with orders. Let $X$ be an $\mathcal{X}$-valued random variable and let $f:\mathcal{X} \to \mathcal{Y}$ be a measurable and non-decreasing function. 
If $X$ satifies the FKG inequality, then so does $f(X)$.
\end{lemma}

We equip the set $\Omega$ introduced in \eqref{eq:def_state_Omega} with the order: 
\begin{equation}
\omega \le \omega'
\qquad 
\text{ if and only if } 
\qquad 
\forall D \in \mathcal{W}^* , \, \omega(D)  \le \omega'(D)\, .
\end{equation}%\forall \omega, \omega' \in \Omega, 
\qquad 

\begin{theorem}\label{th:FKG}
The law $\mathbb{P}$ on $(\Omega, \mathcal{A})$ satifies the FKG inequality.
\end{theorem}

\begin{proof}
It is a general fact that a Poisson point process on a measurable state space with $\sigma$-finite intensity measure satisfies the FKG inequality. See for example Section 20.3 of \cite{Last2017}. 
\end{proof}

Applying Lemma \ref{lem:f(X)_FKG}, we deduce from Theorem \ref{th:FKG} that the following variables satisfy the FKG inequality:
\begin{itemize}
\item the collection $\eta=(\eta_x)_{x\in \mathbb{V}}\in \{0, 1\}^\mathbb{V}$ defined by 
\[\eta_x := 
\begin{cases}
 1 & \text{if } x 
 \text{ is visited by at least one trajectory in }
 \omega \, , \\
0 & \text{else}  \, ,
\end{cases}
\]
the product set $ \{0, 1\}^{\mathbb{V}}$ being equipped with the product order.
In other terms, $\I$ (and therefore also $\V$) satisfies the FKG inequality (with respect to the inclusion order on subsets of $\mathbb{V}$). This is the FKG property shown in \cite{Teixeira2009}.
\item the collection $\eta=(\eta_x)_{x\in \mathbb{V}} \in \N_0^\mathbb{V}$ defined by 
\[\eta_x := 
\text{ number of trajectories in } \omega \text{ visiting } x \, ,
\]
the set $\N_0:=\{0\} \cup \N$ being equipped with the usual order, and the set $\N_0^\mathbb{V}$ with the product order.
\item the collection $\eta=(\eta_x)_{x\in \mathbb{V}} \in \N_0^\mathbb{V}$ defined by 
\[\eta_x := 
\text{ total number of visits to } x
\text{ by trajectories in } \omega\, .
\]
\end{itemize}


% \begin{remark}
%   The process $\sum_{i\ge 1} \delta_{(w_i^*, u_i)}$ of Remark  \ref{rem:full_interlacement_process} also satisfies the FKG inequality, with the order ...
% \end{remark}




\subsection{The \texorpdfstring{$\sigma$}{sigma}-algebra of non-local events}\label{sec:intro_01laws}

We now turn to the main object of this article, which is to study 0-1 laws for random interlacements in our general framework of transient weighted graphs. 

In the case of vertex- or edge-transitive weighted graphs, there is a natural collection of measure-preserving transformations to be considered and one can expect ergodicity, i.e., a 0-1 law for events which are invariant under these transformations. This is the case for the original work of Sznitamn \cite{Sznitman2010}, i.e., for the SRW on $\Z^d, d\ge 3 $, where it was shown that translation-invariant events have probability either 0 or 1.  
In the general case, we rather choose to consider events that are ``non-local'' in the sense that they do not depend on the restriction of the configuration $\omega$ inside any finite region of the graph $\mathcal{G}$. 

\smallskip 

For every finite vertex set $K$, for every biinfinite trajectory $w\in W_K$, we define two relevant times:
\begin{equation}\label{eq:def_ta_K_lambda_k_w}
\begin{aligned}
\tau_K(w) &  = \inf \{n \in \Z : w_n \in K \} \quad \text{ (hitting time of } K),\\
\lambda_K(w) & = \sup \{n\in \Z : w_n \in K\} \quad  \text{ (time of the last visit in } K),
\end{aligned}
\end{equation}
and we break $w$ in three pieces according to these two times: we denote
\begin{itemize}
  \item by $Ins_K(w)$ the finite piece of trajectory of $w$ between times $\tau_K(w)$ and $\lambda_K(w)$, re-indexed so that it starts at time 0: $Ins_K(w):=(w_{\tau_K(w)+n})_{0\le n \le \lambda_K(w)-\tau_K(w)}$, %as in Teixeira but we introduce notations! 
  \item and by $Out_K(w)=(Out_{\rightarrow K}(w) , Out_{K\rightarrow}(w))$ the pair consisiting of the two following semi-infinite pieces of the trajectory $w$:
  \begin{itemize}
  \item $Out_{\rightarrow K}(w)$ the restriction of $w$ up to time $\tau_K(w)$, re-indexed so that it ends at time 0: $Out_{\rightarrow K}(w):=(w_{\tau_K(w)+n})_{n\le 0}$,
  \item and $Out_{K\rightarrow}(w)$ the restriction of $w$ from time $\lambda_K(w)$ on, re-indexed so that it starts at time 0: $Out_{K \rightarrow}(w):=(w_{\lambda_K(w)+n})_{n\ge 0}$.
  \end{itemize} 
\end{itemize}
Furthermore, we define two random variables $\In_K: \omega \mapsto \In_K(\omega)$ and $\Out_K: \omega\mapsto \Out_K(\omega)$ on $(\Omega, \mathcal{A})$ as follows. For any $\omega=\sum_{i\ge 1} \delta_{w_i^*}$ in $\Omega$, we pick for each $i\ge 1$ any representative $w_i$ of $w_i^*$, and define
\begin{equation}
\In_K(\omega)=\sum_{i\in I : w_i^*\in W_K^*} \delta_{Ins_K(w_i)}  \, ,
\end{equation}
and 
\begin{equation}
\Out_K(\omega)=\sum_{i\ge 1: w_i^*\in W_K^*} \delta_{Out_K(w_i)} + \sum_{i\in I : w_i^*\notin W_K^*} \delta_{w_i^*} \, .
\end{equation}
We finally define on $(\Omega, \mathcal{A})$ the sub-$\sigma$-algebra of non-local events:
\begin{equation}
  \mathcal{T}_{RI}= \bigcap_{K \in \mathcal{P}_f(\mathbb{V})} \sigma(\Out_K) \, .
\end{equation}
We will be mainly interested in proving the 0-1 law for the $\sigma$-algebra $\mathcal{T}_{RI}$: it is the property that every event in $\mathcal{T}_{RI}$ has probability either 0 or 1 under $\bbP$. 

\medskip 

This property does not hold in full generality: see Example \ref{ex:simple_not_01} for a simple counter-example. 


\begin{example}\label{ex:simple_not_01}%IN V2: send to section examples.
Consider the Markov associated with the nearest neighbor graph on $\Z$ and the family of conductances:
\begin{equation}
  a_{n, n+1}=2^n, \qquad a_{-n-1, -n}=2^n, \qquad  n\ge 0, 
\end{equation}
i.e., the Markov chain with the transition probabilities $p_{0, 1}=p_{0, -1}=\frac{1}{2}$, and, for $n\ge 1$, $p_{n, n+1}=p_{-n, -n-1}=\frac{2}{3}$ and $p_{n, n-1}=p_{-n, -n+1}=\frac{1}{3}$.
Let us consider the trajectories of $\omega$ that go from $-\infty$ to $\infty$. These trajectories have to pass through 0, so their number is distributed as a Poisson variable with parameter 
\begin{equation}
  \begin{aligned}
  P_0& [X_n \underset{n\to\infty}\longrightarrow -\infty |\tau_0^+=\infty ] e_{\{0\}}(0) P_0[X_n \underset{n\to\infty}\longrightarrow \infty]\\
  = & P_0[X_n \underset{n\to\infty}\longrightarrow -\infty, \tau_0^+=\infty ] a_0 P_0[X_n \underset{n\to\infty}\longrightarrow \infty ].
  \end{aligned}
\end{equation}
This parameter is clearly positive and finite% (it is in fact equal to $1/4$)
, so the number of trajectories in $\omega$ going from $-\infty$ to $\infty$ is a nontrivial random variable. It is however measureable with respect to $\mathcal{T}_{RI}$. 
\end{example}%in fact one could just ask the conductances to be symmetric (and the MC to be transient), so that the two escapes are clearly possible.

More generally, let us give a necessary condition for the 0-1 law to hold.  We consider the sub-$\sigma$-algebra of $(W^+,\mathcal{W}^+)$ consisting of events that are invariant under time-shift. Denoting by 
$\theta: W^+ \to W^+ , \, (w_n)_{n\ge 0}  \mapsto (w_{n+1})_{n\ge 0}$ the time-shift, we set
\begin{equation}
  \mathcal{T}_{MC}:=\left\{ A \in \mathcal{W}^+ : \theta^{-1}(A)=A \right\}.
\end{equation} 
%The events in $\mathcal{T}_{MC}$ describe asymptotic behaviors of the chain $X$ (but independtly of the time origin!!).
For any events $A, B\in\mathcal{T}_{MC}$, let us consider the set of trajectories $w\in W$ such that $(w_{-n})_{n\ge 0}\in A$ and $(w_n)_{n\ge 0}\in B$ and let us denote by ``$A\rightarrow B$'' its image under the projection $\pi^*$. 
We observe that $\omega(A\to B)$, the random variable counting the number of trajectories in $\omega$ that fall in ``$A\rightarrow B$'', is measurable with respect to $\mathcal{T}_{RI}$ and is distributed as a Poisson variable with parameter $\nu(A\rightarrow B)$. 
The announced necessary condition follows: if the 0-1 law for $\mathcal{T}_{RI}$ holds, then
\begin{equation}\label{eq:necessary_cond}
  \forall A, B\in\mathcal{T}_{MC}, \qquad \nu(A\rightarrow B) \in \{0, \infty\}.
\end{equation}


\subsection{Under tail-triviality or tail-atomicity assumption}\label{sec:state_tail_assumptions}


We give a first condition under which the 0-1 law holds.
We say that the Markov chain $X$ is tail-trivial or that the 0-1 law for $\mathcal{T}_{MC}$ is satisfied if for some $x\in\mathbb{V}$, the invariant $\sigma$-algebra $\mathcal{T}_{RI}$ is trivial under $P_x$, i.e., $P_x$ gives probability either 0 or 1 to all event in $\mathcal{T}_{RI}$. Note that, by irreducibilty, requiring this for some $x$ in $\mathbb{V}$ is equivalent to requiring it for all $x$'s in $\mathbb{V}$. 
% See Lemma \ref{lem:equiv_conditions_01_MC_appendix} for equivalent definitions of the tail-triviality of the Markov chain $X$.%888

\begin{theorem}\label{th:tail_triviality}
Assume that the Markov chain $X$ is tail-trivial. Then, under $\bbP$, every event in $\mathcal{T}_{RI}$ has probability either 0 or 1.
\end{theorem}


Next we go beyond the tail-triviality hypothesis. 
Given a probability space $(\mathcal{X}, \mathcal{J}, \mathbb{Q})$, we say that 
\begin{itemize}
\item an event $A\in \mathcal{J}$ is  atomic if $\mathbb{Q}(A)>0$ and any event included in $A$ has probability either equal to that of $\mathbb{Q}(A)$ or to zero,
\item the probability space is purely atomic if there exists an at most countable collection $(A_i)_{i\in I}$ of atomic events partitioning $\mathcal{X}$.
\end{itemize}
We say that $\mathcal{T}_{MC}$ is purely atomic if $\mathcal{T}_{MC}$ is purely atomic with respect to $P_x$ for some $x$ in $\mathbb{V}$.  Note that, by irreducibilty, requiring this for some $x$ in $\mathbb{V}$ is equivalent to requiring it for all $x$'s in $\mathbb{V}$. 
%See Remark \ref{rem:tail_atomicity_Martin} for an equivalent definition in terms of Martin boundary of the Markov chain $X$.  %888
Our next theorem states that, when $\mathcal{T}_{MC}$ is purely atomic, \eqref{eq:necessary_cond} is a sufficient condition for the 0-1 law to hold.
\begin{theorem}\label{th:tail_atomicity}
Assume that $\mathcal{T}_{MC}$ is purely atomic. 
Assume that for every $A, B\in\mathcal{T}_{MC}$, $\nu(A\rightarrow B)$ is either 0 or infinite. 
Then, under $\bbP$, every event in $\mathcal{T}_{RI}$ has probability either 0 or 1.
\end{theorem}


\begin{remark}
  The conditon \eqref{eq:necessary_cond} can also be formulated as: for every $A, B\in \mathcal{T}_{MC}$, as soon as $P_x(A)>0$ and $P_x(B)>0$ for some $x\in \mathbb{V}$ (equivalently, for all $x$'s in $\mathbb{V}$), the quantity $\nu(A\to B)$ must be infinite.

We observe that
$\nu(A\rightarrow B)=\lim_{K\to \mathbb{V}} \nu(A\rightarrow B, W_K^*)$ and that the quantity appearing in this limit can be expressed as
\begin{equation}
\nu(A\rightarrow B, W_K^*)=\sum_{x\in \partial_X K} e_K(x) P_x[ A | \tau_K^+=\infty] P_x[B].
\end{equation}
\end{remark}

\subsection{Quantitative criterion}

We further provide a quantitative criterion for the occurrence of the 0-1 law for $\mathcal{T}_{RI}$.
For every finite vertex set $K$, let us introduce the hinge measure on $K$, the measure on $\mathbb{V}\times \mathbb{V}$ defined by:
\begin{equation}\label{eq:def_hinge_meas}
h_K(x, y)= e_K(x) \times P_x[X_{\lambda_K}= y]\, ,
\end{equation}
which is supported by $\partial_X K \times \partial_X K$ and symmetric (by reversibility of $X$), whose marginals are both equal to the equilibrium measure $e_K$ and whose total mass is equal to $\capa(K)$.

\begin{theorem}\label{th:quantitative_criterion}
  The 0-1 law for $\mathcal{T}_{RI}$ holds if and only if for every site $x\in \mathbb{V}$, for every $\epsilon>0$, as $L\to \mathbb{V}$,
  \begin{equation}
    \sum_{x', y'\in \partial_X L} h_L(x', y') \left(P_{x'}[\tau_x < \infty | X_{\lambda_L}=y']-\epsilon \right)_+ \to 0.
  \end{equation}
\end{theorem}
\begin{remark}
We stress that, due to \eqref{eq:consistency_eq_meas},
\begin{equation}
    \sum_{x', y'\in \partial_X L} h_L(x', y') P_{x'}[\tau_x < \infty | X_{\lambda_L}=y']
    =  \sum_{x'\in L} e_L(x')  P_{x'}[\tau_x<\infty]
    = Cap(\{x\}).
  \end{equation}
The condition in the theorem can therefore be interpreted as the fact that this sum is ``diluted'' enough, that there is not too much mass supported on some couples $(x',y')$.  

Using the bounds $b \ind_{a > 2b} \le (a-b)_+ \le \ind_{a > b}$ for $a \in [0, 1], b \ge 0$, the condition in the theorem can be equivalently formulated as follows: for every site $x\in \mathbb{V}$, for every $\epsilon>0$, as $L\to \mathbb{V}$,
  \begin{equation}
    \sum_{x', y'\in \partial_X L} h_L(x', y') \ind_{P_{x'}[\tau_x < \infty | X_{\lambda_L}=y'] >\epsilon} \to 0.
  \end{equation}
% We also note that the involved probability can be rewritten as follows:
% \begin{equation}
%   P_{x'}[\tau_x < \infty | S_{L_L}=y']= P_{x'}[\tau_x < \tau_{y'} | \tau_{y'}<\infty].
% \end{equation}
% EN FAIT à droite on doit avoir la probab qu'il y ait au moins ue visite à y' après $\tau_x$ (mais ce n'est pas forcément la première)... Donc pas très explicite.
\end{remark}



\subsection{A weak 0-1 law}

Let us also introduce another $\sigma$-algebra, focusing only on the future pieces of trajectory. For every finite vertex set $K$, we define a random variable $\Out_{K\rightarrow} : \omega \mapsto \Out_{K\rightarrow}(\omega)$ on $(\Omega, \mathcal{A})$ as follows. For any $\omega=\sum_{i\ge 1} \delta_{w_i^*}$ in $\Omega$, we pick for each $i\ge 1$ a representative $w_i$ of $w_i^*$ and denote
\begin{equation}
\Out_{K\rightarrow}(\omega)=\sum_{i\in I : w_i^*\in W_K^*} \delta_{Out_{K\rightarrow}(w_i)} + \sum_{i\in I :  w_i^*\notin W_K^*} \delta_{w_i^*} \, .
\end{equation}
The $\sigma$-algebra of non-local future events is defined as
\begin{equation}
  \mathcal{T}_{RI\rightarrow} := \bigcap_{K\in \mathcal{P}_f(\mathbb{V})} \sigma(\Out_{K\rightarrow}) \, .
\end{equation}
This $\sigma$-algebra is a priori not so interesting because it lacks much information. We mention it however, since it has the nice property that it satisfies the 0-1 law in full generality. 
\begin{theorem}\label{th:weak_01law}
Under $\bbP$, every event in $\mathcal{T}_{RI\rightarrow}$ has probability either 0 or 1.
\end{theorem}

\begin{remark}\label{rem:backward_weak}
  By reversibility of the Markov chain $X$, reversing time for all the trajectories in $\omega$ does not change the law of $\omega$ under $\bbP$.
  Thus, we obtain also the 0-1 law for the $\sigma$-algebra of non-local past events, defined by focusing on the past pieces of trajectory $Out_{\rightarrow K}(w)$ instead of the future pieces $Out_{K\rightarrow}(w)$.
\end{remark}



\subsection{A 0-1 law for increasing non-local events}

We now present two other ways to weaken the 0-1 law discussed in Section \ref{sec:intro_01laws}, that may lead to 0-1 laws holding under more general assumptions.

\smallskip

First, one could ask only for the 0-1 law to hold for increasing events belonging to $\mathcal{T}_{RI}$. We observe that condition \eqref{eq:necessary_cond} remains necessary for this modified $0-1$ law to hold. Indeed, if there exist events $A$ and $B$ in $\mathcal{T}_{MC}$ such that $\nu(A\to B) \in (0, \infty)$ then for any positive integer $k$, the event $\{\omega(A\to B ) \ge k\}$ is an increasing event in $\mathcal{T}_{RI}$ which has probability different from 0 and from 1.
So this modified 0-1 law does not hold in full generality, and, in the case where $\mathcal{T}_{MC}$ is tail-atomic, it is in fact equivalent to the 0-1 law for $\mathcal{T}_{RI}$. 

\smallskip

Secondly, we consider the possibility of forgetting the pairings of the pieces of trajectories when defining the trace of configuration $\omega$ outside $K$.
We set, for $\omega=\sum_{i\ge 1} \delta_{\omega_i^*}$, picking for each $i\ge 1$ any representative $w_i$ of $w_i^*$,
\begin{equation}
\widetilde{\Out}_K(\omega):=\sum_{i\ge 1: w_i^*\in W_K^*} \delta_{Out_{\to K}(w_i)} + \sum_{i\ge 1: w_i^*\in W_K^*} \delta_{Out_{K\to}(w_i)} + \sum_{i\in I : w_i^*\notin W_K^*} \delta_{w_i^*} \, 
\end{equation}
and define 
$$\widetilde{\mathcal{T}}_{RI} := \bigcap_{K \in \mathcal{P}_f(\mathbb{V})} \sigma\left(\widetilde{\Out}_K\right). $$
Once again, the 0-1 law for $\widetilde{\mathcal{T}}_{RI}$ does not hold in full generality. For example, we observe that if there exists an event $A\in \mathcal{T}_{MC}$ such that, with $ A^c= W^+\setminus A$, we have $\omega(A\to  A^c)\in (0, \infty)$, then the variable 
$$\omega(A \to  W^+)-\omega( W^+ \to A) = \omega(A \to  A^c)-\omega( A^c \to A)$$ 
is mesurable with respect to $\widetilde{\mathcal{T}}_{RI}$ and non-trivial. 

\medskip

Our final theorem states that, remarkably, combining the two approaches we just explored yields a 0-1 law that holds in full generality.

\begin{theorem}\label{th:increasing_01law}
  Under $\p$, any increasing event in $\widetilde{\mathcal{T}}_{RI} $ has probability either 0 or 1.
\end{theorem}


\begin{remark}
  Of course, by taking complements, also decreasing events in $\widetilde{\mathcal{T}}_{RI} $ have probability 0 or 1 under $\p$.
\end{remark}

 

\section{Preliminaries for the proof of the 0-1 laws}


\subsection{Hinge decomposition of the interlacement process}\label{sec:hinge_decomp}

Given a finite vertex set $K$ and a trajectory $w\in W_K$, recall the definition of the times $\tau_K(w)$ and $\lambda_K(w)$ in \eqref{eq:def_ta_K_lambda_k_w} and define the hinge couple of $w$ as:
\begin{equation}
  Hin_K(w):= (w_{\tau_K(w)}, w_{\lambda_K(w)}) .
\end{equation}
We define the \emph{hinge process} $\Hin_K(\omega)$ associated to $\omega=\sum_{i\ge 1} \delta_{w_i^*}$ by picking for each $i \ge 1$ any representative $w_i$ of $w_i^*$ and setting:
\begin{equation}
  \Hin_K(\omega)=\sum_{i\ge 1 : w_i^*\in W_K^*} \delta_{Hin_K(w_i)}.
\end{equation}
Under $\p$, the random variable $\Hin_K$ is distributed as a Poisson point process with intensity measure the hinge measure $h_K$, introduced in \eqref{eq:def_hinge_meas}. 

Let us describe how to sample $\omega$ under $\p$ with respect to the decomposition along the hinge process. 
To sample the trajectories touching $K$,
\begin{itemize}
\item sample first the hinge process, as a Poisson point process with intensity $h_K$.
\item to each hinge couple $(x, y)$, attach a trajectory $w^*$ which decomposes as $Ins_K(w)$ and $Out_K(w)=(Out_{K\rightarrow} (w), Out_{\rightarrow K}(w))$:
\begin{itemize}
\item $Ins_K(w)$ is sampled according to $P_x[ (X_0, X_1, \dots, X_{\lambda_K})\in \cdot | X_{\lambda_K}=y]$ 
\item $Out_{K\rightarrow} (w)$ is obtained by applying a time reversal to a trajectory sampled according to $P_x[ \cdot |  \tau^+_K=\infty]$,
\item $Out_{\rightarrow K}(w)$ is sampled according to $P_y[ \cdot | \tau^+_K=\infty]$.
\end{itemize} 
\end{itemize}
Complete $\omega$ by sampling the trajectories that do not touch $K$, as usual as a Poisson point process with intensity measure $\nu$ restricted to $W^*\setminus W_K^*$.

\smallskip


We observe that the interlacement process satisfies a spatial Markov property: under $\p$, $\Out_K$ depends on $\In_K$ only through the hinge process $\Hin_K$. Let us thus describe how to go from a finite vertex set K to a larger finite vertex set L: this will be at the heart of our arguments for the 0-1 law.
%So, assume that $\eta = \In_K(\omega)$ is given. To sample $\In_L(\omega)$, independently
To sample $\In_L$ given $\In_K$, independently
\begin{itemize}
  \item extend each path in $\In_K$, say with extremities $(x,y)$, by adding two pieces of trajectory:
\begin{itemize}
  \item on the left, one obtained by applying a time reversal to a trajectory sampled according to $P_x[ (X_0, X_1, \dots, X_{\lambda_{L}})\in \cdot |  \tau_{K}^+=\infty]$,
  \item on the right, one sampled according to $P_y[ (X_0, X_1, \dots, X_{\lambda_{L}})\in \cdot |  \tau_{K}^+=\infty]$,
\end{itemize}
\item sample additional trajectories inside $L$ that do not touch $K$ by sampling biinfinite trajectories under $\nu$ restricted to $W_L^*\setminus W_K^*$ and keeping only the pieces of trajectories inside $L$.
\end{itemize}

%This yields a way to sample $\omega$ progressively along an exhaustion of $\mathbb{V}$.


\subsection{Coupling results}


For $\mu$ and $\nu$ probability measures on an at most countable state space $\mathcal{X}$ (equipped with the $\sigma$-algebra consisting of all its subsets), we denote by $d_{TV}(\mu, \nu)$ their total variation distance, which is defined as
\begin{equation}\label{eq:def_dTV}
  \begin{aligned}
    d_{TV}(\mu, \nu) : & = \sup_{A\subset \mathcal{X}} |\mu(A)- \nu(A)| \\
    & = \sum_{x\in \mathcal{X} : \mu(x)>\nu(x) } \mu(x)- \nu(x) \\
    & = \frac{1}{2} \sum_{x\in \mathcal{X} } |\mu(x)- \nu(x) |\\
    & = \inf_{(X, Y) : X\sim \mu, Y\sim \nu} P(X \neq Y) .
   \end{aligned}
\end{equation}
The above infimum is over couplings $(X, Y)$ of $\mu$ and $\nu$ and it is achieved.

\medskip
Using a classical approach to 0-1 laws, we now establish a sufficient condition for the 0-1 law for $\mathcal{T}_{RI}$ in terms of total variation distances. 
%in terms of total variation distances for the hinge process. 
%We will formulate a criterion for the 0-1 law 
%in terms of total variation distances for the hinge process. 
For any finite vertex set $K$, 
\begin{itemize}
\item we call admissible path inside $K$ any nearest-neighbor path in $\mathcal{G}$ with finite length, with starting time 0 and with extremities belonging to $\partial_X K$.
\item we call admissible configuration inside $K$ any configuration $\eta$ such that $\p(\In_K=\eta)>0$, that is, any configuration of the form $\sum_{i=1}^k \delta_{\gamma_i}$, where $k$ is a finite integer (possibly zero), and for each $i=1, \dots, k$, $\gamma_i$ is an admissible path inside $K$. 
\end{itemize}


\begin{lemma}  \label{lem:suf_cond_0_1_law_distance}
  %The 0-1 law for $\mathcal{T}_{RI}$ holds if and only if for every finite vertex set $K$, for every admissible configuration $\eta$ inside $K$, as $L\to \mathbb{V}$,
  Assume that for every finite vertex set $K$, for every admissible configuration $\eta$ inside $K$, as $L\to \mathbb{V}$,
  \begin{equation}\label{eq:condition_dTV_to0_lem_one_dir}
    d_{TV}\left(\mathbb{P}(\Out_L  \in \cdot | \In_K =\eta),
    \mathbb{P}(\Out_L  \in \cdot | \In_K =0) \right) \to 0.
  \end{equation}
  Then every event in $\mathcal{T}_{RI}$ has probability either 0 or 1 under $\p$.

  Furthermore, the above condition can be restricted to $\eta$'s consisting of a single path.
\end{lemma}
\begin{remark}\label{rem:rewrite_d_TV_cond}
We remark that, due to the Markovian property evoked in Section \ref{sec:hinge_decomp}, we can rewrite the total variation distance in \eqref{eq:condition_dTV_to0_lem_one_dir} as
\begin{equation}
   d_{TV}\left(\mathbb{P}( \Hin_L  \in \cdot | \Hin_K =\rho),
    \mathbb{P}(\Hin_L  \in \cdot | \Hin_K =0) \right),
\end{equation}
  where $\rho$ is the hinge process on $K$ corresponding to $\eta$. 
  So the condition in Lemma \ref{lem:suf_cond_0_1_law_distance} really concerns the dependance of the hinge process $\Hin_L$ on the hinge process $\Hin_K$. 
\end{remark}

\begin{proof}
  We assume that the condition in the lemma holds. Let $E\in \mathcal{T}_{RI}$. Let $K$ be a finite vertex set and $\eta$ an admissible configuration on $K$, and let us consider, under a probability $\mathbf{P}$, a coupling $(\omega, \omega')$ with $\omega$ distributed under  $\mathbb{P}(\cdot | \In_K=\eta)$ and $\omega'$ distributed under $\mathbb{P}(\cdot | \In_K=0)$, such that $\Hin_L(\omega)=\Hin_L(\omega')$ with maximal probability, i.e.,
  \begin{equation}
    \mathbf{P}[\Hin_L(\omega) \neq \Hin_L(\omega')] =  
    d_{TV}\left(\mathbb{P}(\Hin_L \in \cdot | \In_K=\eta),
    \mathbb{P}(\Hin_L \in \cdot | \In_K=0) \right) 
  \end{equation}
  and such that on the event $\{\Hin_L(\omega)=\Hin_L(\omega')\}$, $\omega$ and $\omega'$ coincide outside $L$, in the sense that $\Out_L(\omega)= \Out_L(\omega')$. Since $E$ belongs to $\sigma(Out_L)$, on the latter event $\omega$ realises $E$ if and only if $\omega'$ does. It follows that:
  \begin{equation}
    \begin{aligned}
    & |\p(E|\In_K = \eta)  - \p(E|\In_K  = 0)| \\
    & = |\mathbf{E}[ \ind_{\omega \in E} - \ind_{\omega' \in E} ]| \\
    & \le \mathbf{P}[ \ind_{\omega \in E} \neq  \ind_{\omega' \in E} ] \\
    & \le \mathbf{P}(\Hin_L(\omega ) \neq \Hin_L(\omega'))\\
    & = d_{TV}\left(\mathbb{P}(\Hin_L \in \cdot | \In_K=\eta),
    \mathbb{P}(\Hin_L \in \cdot | \In_K=0) \right) .
    \end{aligned}
  \end{equation}
  Sending $L$ to $\mathbb{V}$, the hypothesis of the lemma yields $\p(E|\In_K= \eta) = \p(E|\In_K=0)$. 
  Since this is true for any admissible $\eta$, we get that, under $\p$, $E$ is independent of the variable $\In_K$.  
  Since the variables $(\In_K, K\in \mathcal{P}_f(\mathbb{V}))$ generate the $\sigma$-algebra $\mathcal{A}$, this implies that, under $\p$, $E$ is independent of $\mathcal{A}$, and in particular of itself, hence $\p(E) \in \{0 , 1\}$. This concludes the law of the 0-1 law.

  \smallskip 

  To justify that one can restrict to $\eta$'s of the form $\eta=\delta_\gamma$, where $\gamma$ is an admissible path inside $\K$, we observe that for every admissible configuration $\eta$ inside $K$ and every admissible path $\gamma$ inside $K$,
\begin{equation}
  \begin{aligned}
  & d_{TV}\left(\mathbb{P}(\Hin_L  \in \cdot | \In_K =\eta+\delta_\gamma),
    \mathbb{P}(\Hin_L  \in \cdot | \In_K =\eta) \right) \\
  & \leq d_{TV}\left(\mathbb{P}(\Hin_L  \in \cdot | \In_K =\delta_\gamma),
    \mathbb{P}(\Hin_L  \in \cdot | \In_K =0) \right) .
  \end{aligned}
\end{equation}
  Indeed, it is just a matter of coupling identically the trajectories corresponding to $\eta$ in both processes. 
  The condition in the lemma derives by reasoning by induction on the number of paths in $\eta$, with the help of the triangular inequality for $d_{TV}$.
\end{proof}

In order to establish the quantitative criterion of Theorem \ref{th:quantitative_criterion}, we will show in Section \ref{sec:necessary_and_sufficient_for_quantitative} that the condition of Lemma \ref{lem:suf_cond_0_1_law_distance} is also necesssary for the 0-1 law for $\mathcal{T}_{RI}$ to hold. This will result from a more general statement that we present now, and that will also be useful (in Section \ref{sec:prelim_tail_assump}) to exploit the tail-triviality or tail-atomicity assumptions of Section \ref{sec:state_tail_assumptions}. 

\smallskip

Let us first recall a useful fact. 
For a (possibly time-inhomogeneous) Markov chain $Z=(Z_n)_{n\ge 0}$ on a countable state space $\mathcal{X}$, 
for any starting points $x$ and $y$ in $\mathcal{X}$, there exists, under a probability measure $\mathbf{P}$, a coupling $(X, Y) = ((X_n)_{n\ge 0} , (Y_n)_{n\ge 0} ) $ of the Markov chain started respectively at $x$ and at $y$ 
such that, almost surely, $X$ and $Y$ coincide after their first meeting time $T:=\inf\{n\ge 0: X_n=Y_n\}$ 
and such that for every $n\ge 0$:
\begin{equation}\label{eq:def_max_coupl}
  \mathbf{P}[T>n]=d_{TV}(P_x[Z_n \in \cdot], P_y[Z_n \in \cdot]).
\end{equation}
This coupling is called the maximal coupling of the Markov chain $Z$, see \cite{Griffeath1975, Pitman1976} or Section 12.4.2 of \cite{Lawler2010}. %it can be constructed using the path decomposition method (probably not a correct formulation, this is only mentioned by Pitman as a comparison)


\begin{lemma}\label{lem:general_equiv_01law}
Consider a (possibly time-inhomogeneous Markov chain) $(Z_n)_{n\ge 0}$ on a discrete countable state space. We fix $x_0\in \mathcal{X}$. The following statements are equivalent:
\begin{itemize}
  \item[(i)] For every $n\ge 0$, for every $x$ and $y$ in $\mathcal{X}$ such that $P_{x_0}(Z_n =x)>0$ and $P_{x_0}(Z_n =y)>0$, we have
  \begin{equation}
  d_{TV}(P_{x_0}(Z_m \in \cdot | Z_n = x), P_{x_0}(Z_m \in \cdot | Z_n = y))  \underset{m\to\infty}{\longrightarrow} 0.
  \end{equation}
  \item[(ii)] Every event in $\mathcal{T}:=\bigcap_{n\ge 0} \sigma(Z_m , m \ge n)$ has probability either 0 or 1 under $P_{x_0}$.
\end{itemize} 
\end{lemma}
\begin{proof} The proof of the implication $(i)\implies (ii)$ follows the same scheme as the proof of Lemma \ref{lem:suf_cond_0_1_law_distance}.
  We assume that $(i)$ holds and consider $E\in\mathcal{T}$. 
  % be a tail event. 
  We fix $n\ge 0$ and $x,y$ in $\mathcal{X}$ such that $P_{x_0}(Z_n =x)>0$ and $P_{x_0}(Z_n =y)>0$. 
  Let us consider now an integer $m \ge n$. We use, under a probability $\mathbf{P}$, a coupling $(X, Y)$ of $P_{x_0}(\cdot | Z_n =x)$ and $P_{x_0}(\cdot | Z_n =y)$ such that $X_m= Y_m$ with maximal probability, i.e.:
  \begin{equation}
    \mathbf{P}[X_m\neq Y_m]  =  d_{TV}(P_{x_0}(Z_m \in \cdot | Z_n = x), P_{x_0}(Z_m \in \cdot | Z_n = y)).
  \end{equation}
  We further require that, on the event $X_m=Y_m$, $X$ and $Y$ coincide after time $m$. Since $E \in \sigma(Z_k, k\ge m)$, on the latter event $X$ realises $E$ if and only if $Y$ does. It follows that:
  \begin{equation}
    \begin{aligned}
    & |P_{x_0}(E | Z_n = x) - P_{x_0}(E | Z_n = y)| \\
    & = |\mathbf{E}[ \ind_{X \in E} - \ind_{Y \in E} ]| \\
    & \le \mathbf{P}[ \ind_{X \in E} \neq  \ind_{Y \in E} ] \\
    & \le \mathbf{P}[X_m \neq Y_m]\\
    & = d_{TV}(P_{x_0}(Z_m \in \cdot | Z_n = x), P_{x_0}(Z_m \in \cdot | Z_n = y))\, .
    \end{aligned}
  \end{equation}
  Sending $m$ to infinity, condition $(i)$ yields that $P_{x_0}(E | Z_n = x) = P_{x_0}(E | Z_n = y) $. 
  We have thus established that, under $P_{x_0}$, $E$ is independent of the variable $Z_n$, for every $n \ge 0$. It follows that, under $P_{x_0}$, $E$ is independent of $\sigma(Z_n , n\ge 0)$, and in particular of itself, hence $P_{x_0}(E) \in \{0 , 1\}$ and $(ii)$ is proved. 

  \smallskip

  Conversely, let us prove the implication $(ii)\implies (i)$. Let us assume that the 0-1 law holds, let us consider $n\ge 0$ and $x,y  \in \mathcal{X}$ such that $P_{x_0}(E | Z_n = x)>0$ and $ P_{x_0}(E | Z_n = y) >0$ and let us show that
  \begin{equation}\label{eq:def_delta}
    \delta :=  \lim_{m\to\infty} d_{TV}(P_{x_0}(Z_m \in \cdot | Z_n = x), P_{x_0}(Z_m \in \cdot | Z_n = y)) 
  \end{equation}
  is zero. Note that the quantity in the limit is non-increasing by elementary arguments.
  Let us consider, under a probability $\mathbf{P}$, $(X,Y)=((X_m)_{m\ge n},(Y_m)_{m\ge n} )$ the maximal coupling of the Markov chain started at time $n$ respectively from $x$ and $y$, in the sense of \eqref{eq:def_max_coupl}. For each $m\ge n$, we consider the subset $E_m$ of $\mathcal{X}$ defined as:
  \[E_m:=\{z \in \mathcal{X}, P_{x_0}(Z_m = z | Z_n = x) > P_{x_0}(Z_m = z | Z_n = y)  \} \,, \]
  which satisfies 
  \begin{equation}\label{eq:PEm_ge_delta}
  P_{x_0}(Z_m \in E_m | Z_n = x) 
   \ge P_{x_0}(Z_m \in E_m | Z_n = y)  
   + \delta 
   \ge \delta,
  \end{equation}
  by the definition of $\delta$ in \eqref{eq:def_delta} and by the definition of the total variation distance in \eqref{eq:def_dTV}. 
  Let us consider the event 
  $E:=\{Z_m \in E_m \text{ for infinitely many } m \text{'s}\}$. 
  %It is a tail event, 
  This event belongs to $\mathcal{T}$, thus $(i)$ yields that $p:=P_{x_0}(E)$ is either 0 or 1. We observe that then also $P_{x_0}( E | Z_n = x)$ and $P_{x_0}(E | Z_n = y)$ are equal to $p$. 
  
  On the one hand, if $p=0$, then $P_{x_0}( Z_m \in E_m \text{ i.o.} |Z_n=x )=0$, which implies that $P_{x_0}(Z_m \in E_m |Z_n=x )  \underset{m\to\infty}\longrightarrow 0$, and then \eqref{eq:PEm_ge_delta} yields that $\delta=0$. 

  On the other hand, if $p=1$, then $\mathbf{P}(Y_m \in E_m \text{ i.o.} ) = P_{x_0}(E|Z_n=y)=1$ so, $\mathbf{P}$-almost surely, there exists an $m\ge n$ (in fact, infinitely many $m$'s) such that $Y_m \in E_m$.
  By maximality of the coupling, % DETAILS NEEDED???, 
  $Y_m\in E_m$ implies almost surely that $T\le m$. Thus, $T$ is almost surely finite and $\delta= \lim_{m\to \infty} \mathbf{P}(T>m)=0$. 
  %BELOW IS REASONING By CONTRADICTION
  % Conversely, let us prove that (ii) implies (i). We reason by contradiction. Let us assume that the 0-1 law holds and that there exist $n\ge 0$ and $x,y  \in \mathcal{X}_n$ such that $P_{x_0}(E | Z_n = x)>0$ and $ P_{x_0}(E | Z_n = y) >0$ and such that
  % \[\delta :=  \lim_{m\to\infty} d_{TV}(P_{x_0}(Z_m \in \cdot | Z_n = x), P_{x_0}(Z_m \in \cdot | Z_n = y)) >0 \]
  % (the limit exists by monotonicity, which follows from elementary arguments). 
  % Let us consider $(Z, Z')$ the maximal coupling of the Markov chain started at time $n$ respectively from $x$ and $y$. For each $m\ge n$, we consider the subset $E_m$ of $\mathcal{X}_m$ defined as:
  % \[E_m:=\{z \in \mathcal{X}_m, P_{x_0}(Z_m = z | Z_n = x) > P_{x_0}(Z_m = z | Z_n = y)  \} \,, \]
  % which satisfies 
  % \begin{equation}\label{eq:PEm_ge_delta}
  % P_{x_0}(Z_m \in E_m | Z_n = x) 
  %  > P_{x_0}(Z_m \in E_m | Z_n = y)  
  %  + \delta 
  %  \ge \delta.
  % \end{equation}
  % by definition of the total variation distance and by  [eq. def of delta].
  % Let us consider the event 
  % $E=\{Z_n \in E_n \text{ for infinitely many } n \text{'s}\}$. 
  % %It is a tail event, 
  % It belongs to $\bigcap_{n\ge 0} \sigma(Z_m , m \ge n)$, thus $(i)$ yields that its probability under $P_{x_0}$, denoted $p$, belongs to $\{0, 1\}$. It follows that $E$ also has probability $p$ under $P_{x_0}(\cdot | Z_n = x) $ and under $P_{x_0}(\cdot  | Z_n = x) $. The fact that $p=0$ would imply that $P_{x_0}[Z_m \in A_m |Z_n=x  ]\underset{m\to\infty}\longrightarrow 0$, which contradicts \eqref{eq:PEm_ge_delta}. So $p=1$ and therefore almost surely, the Markov chain $Z'$ will satisfy $Z'_m \in A_m$ for some $m$ (in fact, infinitely many $m$'s).
  % However, by maximality of the coupling (DETAILS NEEDED), if $Z_m'\in A_m$ then $T\le m$. Thus, $T$ is almost surely finite, which contradicts [eq. def of $\delta$]. 
\end{proof}





We finally provide an estimate on the total variation distance between $Poi(\lambda)$, which denotes the Poisson distribution with parameter $\lambda$ and $Poi(\lambda)+1$, which denotes the law of $X+1$, where $X$ follows $Poi(\lambda)$.
\begin{lemma}\label{lem:dTV_Poisson}
For every $\lambda > 0$, 
\begin{equation}
  d_{TV}(Poi(\lambda), Poi(\lambda)+1) \le \frac{1}{2\sqrt{\lambda}}.
\end{equation}
\end{lemma}


\begin{proof}
With $X\sim Poi(\lambda)$, 
\begin{equation}
  \begin{aligned}
    2 \, d_{TV}(Poi(\lambda), Poi(\lambda)+1)  
    & = \sum_{k\ge 0} \frac{\lambda^k}{k!} e^{-\lambda} \left|\frac{k}{\lambda}-1 \right|\\
    & = E\left[\left|\frac{X}{\lambda} - 1 \right|\right]\\ 
    & \le E\left[\left(\frac{X}{\lambda} - 1 \right)^2 \right]^{\frac{1}{2}}
     = \frac{1}{\sqrt{\lambda}}\, .
  \end{aligned}
\end{equation}
\end{proof}

% \begin{remark}
% In fact, the result we need is only that $d_{TV}(Poi(\lambda), Poi(\lambda)+1) $ goes to 0 as $\lambda \to \infty$. 
% We now give an explicit formula for this distance. Observing that $k\mapsto P(X=k)$ is non-decreasing in $k$ up to $\lfloor \lambda \rfloor$ and non-increasing afterwards, we get that:
% \begin{equation}
%   \begin{aligned}
%     \, d_{TV}(Poi(\lambda), Poi(\lambda)+1)  
%     & = \sum_{k\le \lfloor \lambda \rfloor} P(X=k)-P(X=k-1 )\\
%     & = P(X=\lfloor \lambda \rfloor ).
%   \end{aligned}
% \end{equation}
% Using Stirling's formula, $P(X=\lfloor \lambda \rfloor ) \underset{\lambda  \to\infty}\sim \frac{1}{\sqrt{2\pi \lambda}}$. 
% \end{remark}



\section{Proofs under assumptions on \texorpdfstring{$\mathcal{T}_{MC}$}{TMC}}

In this section, we prove Theorems \ref{th:tail_triviality} and \ref{th:tail_atomicity}.


\subsection{Preliminaries}\label{sec:prelim_tail_assump}


We first show that the total mass of $\nu$ is infinite, so that the total number of trajectories in the interlacement process is almost surely infinite.

\begin{lemma}\label{lem:cap_to_infty}
  As $K\to\mathbb{V}$, $\capa(K)\to \infty$. 
\end{lemma}
\begin{proof}
  We use potential theory for the irreducible and reversible Markov chain $X$ (cf Sections 2.4 and 2.5 in Lyons and Peres). %, in particular pages 32, 36 and 38/39 CITE CORRECTLY). 
  We recall that, for any finite vertex set $K$, a unit flow from $K$ to infinity is a function $\theta:\mathcal{E} \to [0, \infty)$ such that:
  \begin{itemize}
  \item $\forall (x,y)\in \mathcal{E}, \theta(x,y)=-\theta(y,x),$
  \item $\forall x \in \mathbb{V}\setminus K, \sum_{y\sim x } \theta(x,y)=0,$
  \item $\forall x, y \in K$ such that $(x,y)\in \mathcal{E}$, $\theta(x,y)=0$,
  \item $ \sum_{x\in K, y\notin K: x\sim y} \theta(x,y)=1$,
  \end{itemize}
  and that we have
  \begin{equation}
    \frac{2}{\capa(K)} = 
    \inf \left\{
      \sum_{x,y\in \mathbb{V}: x\sim y} \frac{1}{a_{x,y}} \theta(x,y)^2 : 
    \theta \text{ is a unit flow from } K \text{ to infinity}
    \right\}, 
  \end{equation}
  with the convention $1/0=\infty$. %and $\inf \{\infty\} = \infty$. 
  Let us fix a non-empty vertex set $K_0$, for example a singleton. By transience of $X$, $\capa(K_0)>0$, hence there exists a unit flow $\theta_0$ from $K_0$ to infinity such that $\sum_{x,y\in \mathbb{V}: x\sim y} \frac{1}{a_{x,y}} \theta_0(x,y)^2<\infty$. %with finite energy. 
  Then, for every vertex set $K$ containing $K_0$, considering
  \begin{equation}
  \theta_K(x,y):= 
    \begin{cases}
      0 & \text{if } x, y \in K, \\
      \theta_0(x,y)  & \text{else}.
    \end{cases}
  \end{equation}
  which is a unit flow from $K$ to infinity, we get that 
  \begin{equation}
    \frac{2}{\capa(K)}  \le \sum_{(x,y) \in \mathbb{V}^2 \setminus K^2 : x\sim y} \frac{1}{a_{x,y}} \theta_0(x,y)^2
  \end{equation}
  and the latter vanishes as $K\to\mathbb{V}$. The lemma follows.
\end{proof}




When the Markov chain $X$ is tail-trivial, we produce a coupling of $X$ started at two different vertices such that the two copies coincides eventually, up to time-shift.

\begin{lemma}\label{lem:equiv_conditions_01_MC}
  Assume that the Markov chain $X$ satisfies the 0-1 law for $\mathcal{T}_{MC}$. Then, for every $x, y\in \mathbb{V}$, there exists a coupling $X=(X_n)_{n\ge 0} , Y=(Y_n)_{n\ge 0}$ of the Markov chain $X$ started respectively at $x$ and $y$ such that, almost surely, there exist (random) finite times $S$ and $T$ for which
  \begin{equation}
    \forall n\ge 0 , \qquad  X_{S+n}=Y_{T+n}, 
  \end{equation} 
  that is, $X$ and $Y$ coincide eventually, up to time-shift. 
\end{lemma}

\begin{proof}
  Assume that the Markov chain $X$ satisfies the 0-1 law for $\mathcal{T}_{MC}$ and pick $x$ and $y$ in $\mathbb{V}$. 
  Consider an exhaustion $(K_n)_{n\ge 0}$ of $\mathbb{V}$ such that $K_0$ contains $x$ and $y$. 
  The Markov chain $(X_{\lambda_{K_n}})_{n\ge 0}$ is tail-trivial and an application of Lemma \ref{lem:general_equiv_01law} yields:
  \begin{equation}\label{eq:proof_MC_coupling_application_Lemma}
  d_{TV}(P_x[X_{\lambda_{K_n}} \in \cdot], P_y[X_{\lambda_{K_n}} \in \cdot]) \underset{n\to  \infty}\longrightarrow 0
  \end{equation}
  Consider the maximal coupling of the Markov process $(X_{\lambda_{K_n}})_{n\ge 0}$, started from $P_x[X_{\lambda_{K_0}}\in \cdot ]$ and $P_y[X_{\lambda_{K_0}}\in \cdot ]$, by \eqref{eq:proof_MC_coupling_application_Lemma}, the two processes almost surely coincide eventually. Now, given the realisation of this coupling, sample trajectories of the Markov chain $X$, started respectively from $x$ and $y$, conditionally on the exit points prescribed by the coupling. Where the two copies of $(X_{\lambda_{K_n}})_{n\ge 0}$ agree, build the two trajectories of $X$ identically. Then, almost surely, the two trajectories eventually coincide, up to time-shift.
\end{proof}

\begin{corollary}\label{cor:coincide_up_time_shift}
  Assume that the 0-1 law for $\mathcal{T}_{MC}$ is satisfied. 
  Then, for every finite vertex set $K$, for every $x, y\in \partial_X K$, there exists a coupling $(X, Y)$ of $P_x[\cdot | \tau_K^+=\infty]$ and $P_y[\cdot | \tau_K^+=\infty]$ such that, almost surely, there exist (random) finite times $S$ and $T$ for which
  \begin{equation}
    \forall n\ge 0 , \qquad  X_{S+n}=Y_{T+n}, 
  \end{equation} 
  that is, $X$ and $Y$ coincide eventually, up to time-shift. 
\end{corollary}
\begin{proof}
  Consider the Markov chain on $\{x\in \mathbb{V}: P_x(\tau_K=\infty)>0\}$ obtained by conditioning the Markov chain $X$ on the event $\tau_K=\infty$. This Markov chain inherits the tail-triviality of $X$ and Lemma \ref{lem:equiv_conditions_01_MC} applies. 
  Now, to build the coupling in the corollary, just choose any way to couple the first step of $X$ and $Y$, and starting from $X_1$ and $Y_1$ couple $X$ and $Y$ such that they coincide eventually up to time-shift. 
\end{proof}

\subsection{Proof in the case of tail-triviality}\label{sec:proof_tail_triviality}

We now prove Theorem \ref{th:tail_triviality}. The proof follows the idea of \cite{Collin2024}.


\begin{proof}[Proof of Theorem \ref{th:tail_triviality}]
  We assume that the 0-1 law for $\mathcal{T}_{MC}$ holds and we show that the condition of Lemma \ref{lem:criterion_dTV_to_0_law} is satisfied. Let $K$ be a finite vertex set and $\eta$ an admissible configuration inside $K$ consisting of a single path. 
  Let $\gep>0$. 
  Our aim is to find a finite vertex set $M$ and to construct a coupling $(\omega, \omega')$ of $\p( \cdot | \Hin_K =\eta)$ and $\p( \cdot | \Hin_K = 0 )$ such that $\Out_M(\omega)=\Out_M(\omega')$ with probability not smaller than $1-\epsilon$.

  By Lemma \ref{lem:cap_to_infty}, we pick $L$ a vertex set containing $K$ and such that 
  \begin{equation}\label{eq:capa_L-capa_Kgeeps}
  \capa(L)-\capa(K)\ge \epsilon^{-2}.
  \end{equation}
  We observe that the number of trajectories in $\omega$ hitting $L$ is distributed as $Poi(\capa(L)-\capa(K))+1$, while for $\omega'$, this number is distributed as $Poi(\capa(L)-\capa(K))$. Using \eqref{eq:capa_L-capa_Kgeeps} and Lemma \ref{lem:dTV_Poisson}, we construct $\In_L(\omega)$ and $\In_L(\omega')$ such that these numbers coincide with probability not smaller than $1- \epsilon / 2$. 
 
  Now, given $\In_L(\omega)$ and $\In_L(\omega')$ we construct $\omega$ and $\omega'$ as follows.
  If the number of paths in $\In_L(\omega)$ and in $\In_L(\omega')$ are equal,
  \begin{itemize}
    \item we arbitrarily pair each path in $\In_L(\omega)$ with a path in $\In_L(\omega')$, and for each of these pairs, we construct thanks to Corollary \ref{cor:coincide_up_time_shift} the extensions of these paths so that they coincide eventually outside a (random) finite vertex set,
    \item we sample identically the trajectories of $\omega$ and $\omega'$ that do not touch $L$.
  \end{itemize}
 Otherwise, we construct $\Out_L(\omega)$ and $\Out_L(\omega')$ independently of each other.

 For this coupling, we have, for a large enough finite vertex set $M$, that $\omega$ and $\omega'$ coincide outside $M$ with probability not smaller than $1-\epsilon$. The proof is complete.
\end{proof}
  
% \begin{remark}DANS LA V2 et REPHRASE
%   In this proof, we have constructed a coupling which suceeds with probability at least $1-\epsilon$. 
%   We can construct a coupling of two configurations with different prescriptions inside K, such that they coincide almost surely outside some random set. 
% \end{remark}


\subsection{Proof in the case of tail-atomicity}\label{sec:proof_tail_atomicity}


We now prove Theorem \ref{th:tail_atomicity}. The proof relies on a decomposition of the interlacement process $\omega$ as a sum of (independent) Poisson point processes similar to the interlacement process under the assumption of tail-triviality. 

\begin{proof}[Proof of Theorem \ref{th:tail_atomicity}]
We assume that $\mathcal{T}_{MC}$ is purely atomic, so we consider a collection $\mathcal{C}$ of events in $\mathcal{T}_{MC}$ that, for arbitrary $x$ in $\mathbb{V}$, are atoms of $(W^+, \mathcal{T}_{MC}, P_x)$ and form a partition of $W^+$. 

For every $A, B \in \mathcal{C}$, let us denote by $\omega_{|A\to B}$ the restriction of $\omega$ to trajectories belonging to the event $``A\to B''$ (defined in the end of Section \ref{sec:intro_01laws}). 
Naturally, the random variable $\omega_{|A\to B}$ is a Poisson point process with intensity measure $\nu \ind_{A\to B}$. 
Let us now describe $\omega_{|A\to B}$ similarly as we did in Section \ref{sec:def_interl_process} for the complete interlacement process $\omega$. 
For every finite vertex set $K$, let
\begin{equation}
  e_{K, A\to B}(x):=
  \begin{cases}
    a_x P_x[A, \tau_K^+=\infty]  P_x[B] & \text{if } x\in  K,\\ 
    0 & \text{otherwise},
  \end{cases}
\end{equation}
\begin{equation}
  {\capa}_{A\to B}(K):=e_{K, A\to B}(\mathbb{V})= \sum_{x\in K} a_x P_x[A, \tau_K^+=\infty]  P_x[B],
\end{equation}
and
\begin{equation}
  \harm_{K, A\to B}(\cdot):= \frac{e_{K, A\to B}(\cdot)}{{\capa}_{A\to B}(K)}\, .
\end{equation}

To sample the trajectories of $\omega_{|A\to B}$ touching $K$:
\begin{itemize}
\item sample a number $\xi$ as a Poisson variable of parameter $\capa_{A\to B}(K)$,
\item independently for each $i=1,  \dots, \xi$, sample a point $x_i$ under $\harm_{K, A\to B}$,
\item independently for each $i=1, \dots, \xi$, attach to $x_i$ a trajectory that enters in $K$ through $x$, at time 0:
\begin{itemize}
  \item for the negative indices, sample a trajectory under $P_x[ \cdot |A, \tau_K^+=\infty]$ and apply the function $n\mapsto -n$ to the indices,
  \item for the positive indices, sample a trajectory under $P_x[ \cdot |B ]$.
\end{itemize}
\end{itemize}

We can also formulate this in terms of hinge process. Define 
\begin{equation}
h_{K, A\to B}(x, y)= e_{K, A\to B}(x) \times P_x[X_{\lambda_K}= y | B ] = a_x P_x[A, \tau_K^+=\infty]  P_x[B, X_{\lambda_K}= y] .
\end{equation}
To sample the trajectories of $\omega_{|A\to B}$ touching $K$,
\begin{itemize}
\item sample first the hinge process, as a Poisson point process with intensity $h_{K, A\to B}$.
\item to each hinge couple $(x, y)$, attach a trajectory $w^*$ which decomposes as $In_K(w)$ and $Out_K(w)=(Out_{K\rightarrow} (w), Out_{\rightarrow K}(w))$:
\begin{itemize}
\item $In_K(w)$ is sampled according to $P_x[ (X_0, X_1, \dots, X_{\lambda_K})\in \cdot | B, X_{\lambda_K}=y]$ 
\item $Out_{K\rightarrow} (w)$ is obtained by applying a time reversal to a trajectory sampled according to $P_x[ \cdot | A, \tau^+_K=\infty]$,
\item $Out_{\rightarrow K}(w)$ is sampled according to $P_y[ \cdot | B, \tau^+_K=\infty]$.
\end{itemize} 
\end{itemize}
For $K\subset L$, sampling $\In_L(\omega_{A\to B})$ given $\In_K(\omega_{A\to B})$ is also very similar to what we presented in Section \ref{sec:hinge_decomp}. 
The crucial observation is that the past pieces of trajectory are sampled under the Markov chain $X$ conditioned on $A$ while the future pieces are sampled under the Markov chain $X$ conditioned on $B$. Since $A$ and $B$ are atoms of $\mathcal{T}_{MC}$, we note that the two Markov chains obtained after these conditionings are tail-trivial, so that Corollary \ref{cor:coincide_up_time_shift} applies.

\smallskip

We now conclude the proof of Theorem \ref{th:tail_atomicity}. 
Consider $A, B \in \mathcal{C}$. Assume that $\nu(A\to B)$ is infinite. Then we reproduce the proof of Theorem \ref{th:tail_triviality} in Section \ref{sec:proof_tail_triviality}, to show that $\omega_{|A\to B}$ satisfies the 0-1 law under $\mathbb{P}$: for a given vertex set $K$ and a given admissible finite path inside $K$, we want to construct a coupling of two copies of $\omega_{|A\to B}$, one conditioned on having exactly this path inside $K$ and one on having no trajectories touching $K$, such that they coincide outside a large vertex set with large probability. We proceed as follows:
\begin{itemize}
\item since $\capa_{A\to B}(K)\underset{K \to \mathbb{V} }\longrightarrow \nu(A\to B) = \infty$, we find a large enough finite vertex set $L$ and make the number of trajectories hitting $L$ in both copies of $\omega_{|A\to B}$ coincide with large probability,
\item we arbitrarily pair the trajectories touching $L$ between both copies of $\omega_{|A\to B}$ and use Corollary \ref{cor:coincide_up_time_shift} for the Markov chain $X$ conditioned respectively on $A$ and on $B$ to make them coincide outside a large vertex set $M$, with large probability,
\item we sample the trajectories that do not touch $L$ identically in both copies of $\omega_{|A\to B}$.
\end{itemize}
We deduce that the variable $\omega_{|A\to B} $ satisfies the 0-1 law under $\mathbb{P}$. 


Now, if $\nu(A\to B)$ is infinite for every $A, B \in \mathcal{C}$, then all the variables $(\omega_{|A\to B})_{A,B  \in \mathcal{C} }$ satisfy the 0-1 law under $\mathbb{P}$, and so does their sum $\omega$.
This concludes the proof of Theorem \ref{th:tail_atomicity}.
\end{proof}

\begin{remark}
We stress that the variables $(\omega_{|A\to B})_{A,B  \in \mathcal{C} }$ form an independent family, although this does not play a role in the proof. 
\end{remark}




\section{Quantitative approach}\label{sec:quantitative}

\subsection{Preliminaries}

\subsubsection{Necessary and sufficient condition}\label{sec:necessary_and_sufficient_for_quantitative}

For any finite vertex set $K$, 
\begin{itemize}
\item we call admissible hinge couple on $K$ any couple $(x,y)\in \partial_X K \times \partial_X K $;
\item we call admissible hinge configuration on $K$ any configuration $\rho$ such that $\p(\Hin_K=\rho)>0$, that is, any configuration of the form $\sum_{i=1}^k \delta_{(x_i,y_i)}$, where $k$ is a finite integer (possibly zero), and for each $i=1, \dots, k$, $(x_i,y_i)$ is an admissible hinge couple on $K$. 
\end{itemize}

\begin{lemma}\label{lem:criterion_dTV_to_0_law}
  The 0-1 law for $\mathcal{T}_{RI}$ holds if and only if for every finite vertex set $K$, for every admissible hinge configuration $\rho$ on $K$, as $L\to \mathbb{V}$,
  \begin{equation}\label{eq:criterion_d_TV_cond}
    d_{TV}\left(\mathbb{P}(\Hin_L \in \cdot | \Hin_K=\rho),
    \mathbb{P}(\Hin_L \in \cdot | \Hin_K=0) \right) \to 0.
  \end{equation}
  Furthermore, the condition can be restricted to $\rho$'s that consist of a single hinge couple.
\end{lemma}

\begin{proof}
  We have already shown that \eqref{eq:criterion_d_TV_cond} implies the 0-1 law for $\mathcal{T}_{RI}$: this is Lemma \ref{lem:suf_cond_0_1_law_distance}, together with Remark \ref{rem:rewrite_d_TV_cond}. 
  
  To show the converse implication, we use Lemma \ref{lem:general_equiv_01law}. Assume that the 0-1 law for $\mathcal{T}_{RI}$ holds and consider a finite vertex set $K$, an admissible hinge configuration $\rho$ on $K$, and an exhaustion $(L_n)_{n\ge 0}$ of $\mathbb{V}$. Assume without loss of generality that $L_0=\emptyset$, $L_1=K$. Apply Lemma \ref{lem:general_equiv_01law} to the time-inhomogeneous Markov chain $(\Hin_{K_n})_{n\ge 0}$ to obtain:
  \begin{equation}
    d_{TV}\left(\mathbb{P}(\Hin_{L_n} \in \cdot | \Hin_K=\rho),
    \mathbb{P}(\Hin_{L_n} \in \cdot | \Hin_K=0) \right) 
    \underset{n\to\infty}\longrightarrow 0.
  \end{equation}
  This concludes the proof of the lemma.
\end{proof}

% DANS LA V2 
%REMARQUE A ECRIRE, MAIS SANS DOUTE PAS ICI ??? : using a maximal coupling, the condition implies that we can construct a coupling of the processes under the 2 conditionings such that they coincide outside a (random) finite vertex set.




\subsubsection{Distance between Poisson point processes}

If $\nu$ is a finite measure on a finite set $\mathcal{X}$, let us denote by $PPP(\nu)$ the law of a Poisson point process on $\mathcal{X}$ with intensity measure $\nu$. If furthermore, $\pi$ is a probability measure on $\mathcal{X}$, let us denote by $PPP(\nu) \oplus \pi$ the law of the superposition of this same Poisson point process and a random point in $\mathcal{X}$ sampled independently under $\pi$. 

\begin{lemma}\label{lem:general_cond_dTVto0_for_PPP}
For every $n\ge 1$ we consider a finite set $\mathcal{X}_n$, and two finite measures $\nu_n$ and $\pi_n$ on $\mathcal{X}_n$, $\pi_n$ being a probability measure. Then, 
\begin{equation}
  d_{TV}(PPP(\nu_n) \oplus \pi_n  , PPP(\nu_n) ) \underset{n\to\infty}{\longrightarrow} 0
\end{equation}
if and only if, for every $\epsilon>0$, 
\begin{equation}
  \sum_{x\in \mathcal{X}_n} 
  \left(\pi_n(x)-\epsilon \nu_n(x)\right)_+  
  \underset{n\to\infty}{\longrightarrow} 0\, .
\end{equation}
\end{lemma}


% \begin{remark}
% Using the bounds $\frac{a}{2} \ind_{a > 2b} \le (a-b)_+ \le a \ind_{a > b}$ for $a,b \ge 0$, 
% the condition in the lemma is equivalent to the fact that for every $\epsilon>0$, 
% \begin{equation}
%   \sum_{x\in \mathcal{X}_n} 
%   \pi_n(x) \ind_{\pi_n(x) > \epsilon \nu_n(x)} 
%   \underset{n\to\infty}{\longrightarrow} 0\, .
% \end{equation}
% \end{remark}


\begin{proof}
  We consider a finite state space $\mathcal{X}$, and two measures $\nu$ and $\pi$ on $\mathcal{X}$, $\pi$ being a probability measure. 

  Let us first show, for every  $\epsilon>0$ and $a>0$, that if $\sum_{x\in \mathcal{X}} 
  \left(\pi(x)-\epsilon \nu(x)\right)_+  \le a$, then  
\begin{equation}\label{eq:dTV_PPP_upp}
  d_{TV}( PPP(\nu) \oplus \pi, PPP(\nu)) \le \frac{1}{2}\sqrt{\frac{\epsilon}{1-a}} + a ,
\end{equation}
To prove this, we construct $PPP(\nu)$ using an auxiliairy Poisson point process. We consider $\sum_{i\ge 1} \delta{(x_i, u_i)}$ a Poisson point process on $\mathcal{X} \times [0, \infty)$ with intensity measure $\rho=\mathbf{c} \otimes du$, where $\mathbf{c}:=\sum_{x\in \mathcal{X}} \delta_x$ is the counting measure on $\mathcal{X}$ and $du$ is the Lebesgue measure on $(0, \infty)$. 
Then 
$\omega_\nu:=\sum_{i\ge 1: u_i\le \nu(x_i)} \delta_{x_i}$ 
realises $PPP(\nu)$. 
Furthermore, let us consider the domain
$$D=\left\{(x,u)\in \mathcal{X}\times (0, \infty): u\le \frac{\pi(x)}{\epsilon} \wedge \nu(x)\right\}.$$
Let us tilt $\omega_\nu$ only by changing the law of the number of points falling in $D$: let us make it follow the law of $Poi(\rho(D))+1$ instead of the law $Poi(\rho(D))$. 
Then, we obtain a point process which is the superposition of a $PPP(\nu)$ and a point sampled under the renormalisation of the measure $\frac{\pi(\cdot)}{\epsilon} \wedge \nu(\cdot)$. This last distribution is at total-variation distance not larger than $a$ from $\pi$.
Thus:
\begin{equation}
  d_{TV}( PPP(\nu) \oplus \pi, PPP(\nu)) \le d_{TV}(Poi(\rho(D)), Poi(\rho(D))+1)+ a ,
\end{equation}
We conclude the proof of \eqref{eq:dTV_PPP_upp} by using Lemma \ref{lem:dTV_Poisson} and observing that $\rho(D) \ge \frac{1-a}{\epsilon}$.
% \begin{equation}
%   \rho(D) 
%   = \sum_{x\in \mathcal{X}} \frac{\pi(x)}{\epsilon} \wedge \nu(x) 
%   = \frac{1}{\epsilon} \left(1 - \sum_{x\in \mathcal{X}} (\pi(x) -\epsilon \nu(x))_+ \right)
%   \ge \frac{1-a}{\epsilon} 
% \end{equation}

\medskip

Let us furthermore show for every $\epsilon >0$ and $a>0$ the existence of a positive real $D_{\epsilon, a}$ such that if 
$\sum_{x\in \mathcal{X}} \left(\pi(x)-\epsilon \nu(x)\right)_+  \ge a$, then 
  \begin{equation}\label{eq:dTV_PPP_low}
  d_{TV}( PPP(\nu) \oplus \pi, PPP(\nu)) \ge D_{\epsilon, a}.
\end{equation}
Indeed, let us consider $A:=\{ x \in \mathcal{X}: \pi(x)  > \epsilon \nu(x)\}$. We have by assumption $\pi(A) \ge \epsilon \nu(A) +a$. Considering the number of points falling in $A$, we get the lower bound:
$$d_{TV}( PPP(\nu) \oplus \pi, PPP(\nu))
\ge d_{TV}( Poi(\nu(A))+Ber(\pi(A)), Poi(\nu(A))) 
\ge D_{\epsilon, a} \, , $$
where we denote by $Poi(\lambda)+Ber(p)$ the law of the independent sum of a Poisson variable of parameter $\lambda$ and a Bernoulli variable of parameter $p$ and we set 
$$D_{\gep, a}  := 
\inf_{\lambda \ge 0, \,  p\in [0, 1]:\, \epsilon \lambda +a\le p} 
d_{TV}(Poi(\lambda)+Ber(p), Poi(\lambda)).$$
We observe that the function 
$$ (\lambda,p) \mapsto  d_{TV}(Poi(\lambda)+Ber(p), Poi(\lambda))$$
defined on the compact   
$\{(\lambda, p): \lambda \ge 0, p\in [0, 1], \epsilon \lambda+a\le p \}$ 
is continuous and positive, hence $D_{\gep, a}>0$. 
% Furthermore, as $x\to\infty$, it converges towards 1, using the law of large numbers ($Poi(x)/x\to 1$ while $Poi((1+\frac{1}{a}) x +\gep)/x \to 1+\frac{1}{a}$).
Thus, \eqref{eq:dTV_PPP_low} is established and the lemma follows from \eqref{eq:dTV_PPP_upp} and \eqref{eq:dTV_PPP_low}.
\end{proof}




\subsection{Proof of the quantitative criterion}\label{sec:proof_quantitative}

We now prove Theorem \ref{th:quantitative_criterion}.

\begin{proof}[Proof of Theorem \ref{th:quantitative_criterion}]
Assume that the 0-1 law for $\mathcal{T}_{RI}$ holds under $\mathbb{P}$. Consider a vertex $x$ in $\mathbb{V}$, set $K=\{x\}$ and consider the hinge couple $(x,x)$ on $K$. Applying Lemma \ref{lem:criterion_dTV_to_0_law}, we get the convergence towards zero, as $L\to\mathbb{V}$, of the total variation distance between 
\begin{itemize}
  \item the law 
  $\mathbb{P}(\Hin_L \in \cdot | \Hin_K=0)$, that is, a Poisson point process with intensity measure 
  $h_{L}^{x}(x',y') := h_L(x',y') P_{x'}(\tau_x=\infty|X_{\lambda_L}=y')$,
  \item and the law 
  $\mathbb{P}(\Hin_L \in \cdot | \Hin_K=\delta_{(x,x)})$, that is, the law of the superposition of the above Poisson point process and an independent point sampled under the probability 
  $p_{L}^{x}(x', y'):= P_x(X_{\lambda_L}= x'| \tau_x^+=\infty) P_x(X_{\lambda_L}= y'| \tau_x^+=\infty)$.
\end{itemize}
Applying Lemma \ref{lem:general_cond_dTVto0_for_PPP} for any exhaustion $(L_n)_{n\ge 0}$ of $\mathbb{V}$, this implies that for every $\epsilon>0$, 
\begin{equation}
  \sum_{x',y' \in \partial_X L} 
  \left(p_{L}^{ x}(x', y') 
  - \epsilon h_{L}^{x}(x',y') \right)_+  \underset{L \to \mathbb{V}}\longrightarrow 0.
\end{equation}
Bounding $h_{L}^{x}$ by $h_L$, this implies that  
\begin{equation}
  \sum_{x',y'\in  \partial_X L} 
  \left(p_{L}^{x}(x', y')  - \epsilon h_L(x', y')   \right)_+  
  \underset{L \to \mathbb{V}}{\longrightarrow} 0.
\end{equation}
Now, by reversibility:
\begin{equation}
  a_x  P_x[X_{\lambda_L}=x', \tau_x^+=\infty]    P_x[X_{\lambda_L}=y']  = a_{x'} P_{x'}[\tau_L=\infty]  P_{x'}[\tau_x < \infty, X_{  \lambda_L}=y'].
\end{equation}
We also note that $P_x[X_{\lambda_L} = z,\tau_x^+=\infty]= P_x[\tau_x^+=\infty] P_x[X_{\lambda_L} = z]$, for every $z\in \mathbb{V}$.
It follows that:
\begin{equation}\label{eq:rewrite_proba_induced_hinges_singleton}
 p_{L}^x(x', y') =  \frac{h_L(x', y')}{a_x P_x(\tau_x^+=\infty) } 
  P_{x'}[\tau_x < \infty | X_{\lambda_L}=y'].
\end{equation}
Recalling that $x$ is fixed, and playing on $\epsilon$, we deduce that for every $\epsilon>0$
\begin{equation}
  \sum_{x',y' \in \partial_X L} 
  h_L(x', y')
  \left(P_{x'}[\tau_x < \infty | X_{\lambda_L}=y']
  - \epsilon  \right)_+  \underset{L \to \mathbb{V}}{\longrightarrow} 0.
\end{equation}
This is the condition of Theorem \ref{th:quantitative_criterion}.

\medskip 

Conversely, let us assume that the condition of Theorem \ref{th:quantitative_criterion} holds, and let us pick a finite vertex set $K$, and $(x,y)$ an admissible hinge couple on $K$. Following Lemma \ref{lem:general_cond_dTVto0_for_PPP}, our aim is to show the convergence towards zero, as $L\to\mathbb{V}$, of the total variation distance between
\begin{itemize}
  \item the law 
  $\mathbb{P}(\Hin_L \in \cdot | \Hin_K=0)$, that is, a Poisson point process with intensity measure 
  $h_{L}^{K}(x',y') := h_L(x',y') P_{x'}(\tau_K=\infty|X_{\lambda_L}=y')$
  \item and the law 
  $\mathbb{P}(\Hin_L \in \cdot | \Hin_K=\delta_{(x,y)})$, that is, the law of the superposition of the above Poisson point process and an independent point sampled under the probability 
  $p_{L}^{K, x,y} (x', y'):= P_x(X_{\lambda_L}= x'| \tau_K^+=\infty) P_y(X_{\lambda_L}= y'| \tau_K^+=\infty)$.
\end{itemize}
By Lemma \ref{lem:general_cond_dTVto0_for_PPP}, it is enough to establish that for every $\epsilon>0$, 
\begin{equation}\label{eq:sum_pLKxy-epshLK_to0}
  \sum_{x',y' \in \partial_X L} 
  \left(p_L^{K, x, y}(x',y') 
  - \epsilon h_L^{K}(x',y') \right)_+  \underset{L \to \mathbb{V}}{\longrightarrow} 0.
\end{equation}
By reversibility:
\begin{equation}
  \begin{aligned}
  a_x &  \p_x[X_{\lambda_L}=x', \tau_K^+=\infty]   \p_x[X_{\lambda_K}=y] \p_y[X_{\lambda_L}=y' | \tau_K^+=\infty] \\
   =  & a_{x'} \p_{x'}[\tau_L^+=\infty]
   \p_{x'}[\tau_K < \infty, X_{\tau_K} = x, X_{\lambda_K}=y, X_{\lambda_L}=y']
  \end{aligned}
\end{equation}
hence
\begin{equation}
  p_L^{K, x, y}(x',y')= \frac {h_L(x', y')}{ h_K(x, y)} 
  \p_{x'}[\tau_K < \infty, X_{\tau_K} = x, X_{\lambda_K}=y| X_{\lambda_L}=y'],
\end{equation}%\label{eq:rewrite_proba_induced_hinges}
and therefore, playing on $\epsilon$, \eqref{eq:sum_pLKxy-epshLK_to0} is equivalent to the fact that for every $\epsilon>0$, 
\begin{equation}
  \begin{aligned}
    \sum_{x', y'\in \partial_X L} &  h_L(x', y') \\
    &\times \big(  P_{x'}[\tau_K<\infty, X_{\tau_K}=x, X_{\lambda_K}=y |X_{\lambda_L}=y']  -\epsilon P_{x'}(\tau_K=\infty|X_{\lambda_L}=y') \big)_+ 
  \end{aligned}
  \end{equation}
vanishes as $L \to \mathbb{V}$. Bounding $P_{x'}[\tau_K<\infty, X_{\tau_K}=x, X_{\lambda_K}=y  |X_{\lambda_L}=y']$ by $P_{x'}[\tau_K<\infty|X_{\lambda_L}=y']$, the latter follows from the fact that for every $\epsilon>0$,
\begin{equation}
    \sum_{x',y' \in \partial_X L} h_L(x', y') \left((1+\epsilon) P_{x'}[\tau_K<\infty | X_{\lambda_L}=y']-\epsilon \right)_+ \underset{L \to \mathbb{V}}{\longrightarrow} 0,
  \end{equation}
which, finally, can be derived elementary from the condition of Theorem \ref{th:quantitative_criterion} by using the bound:
\begin{equation}
  P_{x'}[\tau_K<\infty|X_{\lambda_L}=y'] \le \sum_{z\in K}  P_{x'}[\tau_z<\infty|X_{\lambda_L}=y']
\end{equation}
and playing on $\epsilon$.
\end{proof}





\subsection{Proof of the weak 0-1 law}


This section is devoted to showing Theorem \ref{th:weak_01law}, that is, the 0-1 law for events that depend, for every finite vertex set $K$, only on the future pieces of trajectory outside $K$. 

Let us observe that in the case of tail-atomicity, one could reproduce our proof (in Section \ref{sec:proof_tail_atomicity}) to show that this 0-1 law holds as soon as $\nu(W^+ \to A)\in \{0, \infty\}$ for every $A\in \mathcal{T}_{MC}$. 
And in fact, for $A$ an atom of $\mathcal{T}_{MC}$ under $P_x$ (for arbitrary $x$), applying Lemma \ref{lem:cap_to_infty} to the Markov chain $X$ conditioned on $A$ yields that $\nu (A \to A)=\lim_{K\to\mathbb{V}} \capa_{A\to A}(K)= \infty$ and a fortiori $\nu(W^+ \to A)=\infty$. 

\smallskip

Since we aim to prove Theorem \ref{th:weak_01law} in full generality, we will rather rely on the quantitative approach of the present Section \ref{sec:quantitative}. 
We will proceed in two steps. 
We will first prove that a condition similar to that of Theorem \ref{th:quantitative_criterion} holds in full generality. Then, we will prove Theorem \ref{th:weak_01law}, using the same scheme of proof as in Section \ref{sec:proof_quantitative}.

\begin{lemma}\label{lem:quantitative_weak}
  For every site $x\in \mathbb{V}$, for every $\epsilon>0$, as $L \to \mathbb{V}$,
  \begin{equation}
    \sum_{x'\in L} e_L(x') \left(P_{x'}[\tau_x < \infty]-\epsilon \right)_+ \to 0.
  \end{equation}
\end{lemma}

\begin{remark} \label{rem:rewrite_quantitative_weak}
  We stress that, due to \eqref{eq:consistency_eq_meas},
\begin{equation}
    \sum_{x'\in L} e_L(x') P_{x'}[\tau_x < \infty]= Cap(\{x\}).
  \end{equation}
The content of the lemma can therefore be interpreted as the fact that this sum is ``diluted'' enough, that there is not too much mass supported on some vertices $x$.  

Using the bounds $b \ind_{a > 2b} \le (a-b)_+ \le \ind_{a > b}$ for $a \in [0, 1], b \ge 0$, the condition in the theorem can be equivalently formulated as follows: for every site $x\in \mathbb{V}$, for every $\epsilon>0$, as $L\to \mathbb{V}$,
  \begin{equation}
    \sum_{x'\in  L} e_L(x') \ind_{P_{x'}[\tau_x < \infty ] >\epsilon} = e_L(\{x'\in L : P_{x'}[\tau_x < \infty ] >\epsilon \} )\to 0.
  \end{equation}
\end{remark}


\begin{proof}[Proof of Lemma \ref{lem:quantitative_weak}] 
  Let us fix $x\in \mathbb{V}$ and $\epsilon >0$. We consider the set 
  \[D:=\{y \in \mathbb{V} : P_y[\tau_x<\infty]> \epsilon\}.\]
  As observed in Remark \ref{rem:rewrite_quantitative_weak}, our aim is to show that $e_L(D)$ vanishes as $L\to\mathbb{V}$.
  
  Of course, if $D$ is finite this is immediate: as soon as $L$ contains $\overline{D}: =\{x\in \mathbb{V}:\exists y \in D, x\sim y\}$, we have $\partial_X L \cap D =\emptyset$ and thus $e_L(D)=0$. But $D$ can be infinite. Think for example of a simple random walk on $\mathbb{Z}$ with a positive drift and $x=0$: in this example, $P_y[\tau_x<\infty]=1$ for every $y\le 0$.

  Let us show that $D$ is however necessarily transient, i.e., that for every $z\in \mathbb{V}$, 
  $P_z[X_n \text{ visits } D \text{ infinitely many times}]=0$. 
  Our proof relies on the fact that $X$ visits $x$ almost surely finitely many times, due to transience.

  For any $y$ in $D$, we pick an integer $t_y$ such that $P_y[\tau_x <t_y]\ge \frac{\epsilon}{2}$. Then we decompose the trajectory of the Markov chain $X$ starting from $z$ in disjoint excursions:
  \begin{itemize}
  \item starting from $z$, we wait for the first visit to $D$, say at time $N_1$ and at site $Y_1$; the first excursion is the piece of trajectory from time $N_1$ to time $N_1+t_{Y_1}-1$;
  \item we reiterate: for $n\ge 1$, after the $n$-th excursion, we wait for the first visit to $D$, say at time $N_{n+1}$ and site $Y_{n+1}$, then we define the $(n+1)$-th excursion to be the piece of trajectory from time $N_{n+1}$ to time $N_{n+1}+t_{Y_{n+1}}-1$. 
  \end{itemize}
  There may be only finitely many excursions, meaning that at some point the trajectory stops visiting $D$; this is actually what we want to show happens almost surely.
  We observe that every excursion has a probability at least $\frac{\epsilon}{2}$ of visiting $x$. Thus, the expected number of excursions is not larger than $\frac{2}{\epsilon}$ times the expected number of visits to $x$, which is finite by transience of $X$. We deduce that the number of excursions is almost surely finite, hence $D$ is transient.

  To conclude, we observe, using the definition of $D$ and reversibility that
  \begin{align*}
  e_L(D) & =    \sum_{y\in D} a_y P_y[\tau_L^+=\infty]\\
    & \le \frac{1}{\epsilon} \sum_{y\in D} a_y P_y[\tau_L^+=\infty]  P_y[\tau_x<\infty] \\
    & = \frac{1}{\epsilon} \sum_{y\in D} a_x P_x[X_{\lambda_L}=y]
     = \frac{a_x}{\epsilon} P_x[X_{\lambda_L}\in D].
  \end{align*}
  The transience of $D$ implies that $P_x[X_{\lambda_L}\in D]$ vanishes as $L\to\mathbb{V}$ and the proof is complete.
\end{proof}
  
  


\begin{proof}[Proof of Theorem \ref{th:weak_01law}]
  For convenience, we prefer to show the 0-1 law for the $\sigma$-algebra of non-local backward events, mentioned in Remark \ref{rem:backward_weak} and defined by focusing on the past pieces of trajectory $Out_{\rightarrow K}(w)$ instead of the future pieces $Out_{K\rightarrow}(w)$. As justified in that remark, these two ``weak'' 0-1 laws are equivalent.%CE N'EST pas vraiment nécessaire...

  Analogously to Lemma \ref{lem:criterion_dTV_to_0_law} (or Lemma \ref{lem:suf_cond_0_1_law_distance}), it is enough to show that for every finite vertex set $K$, for every $x\in \partial_X K$, as $L\to\mathbb{V}$,
  \begin{equation}\label{eq:criterion_d_TV_cond_bis}
    d_{TV}\left(\mathbb{P}(\Hin_{\to L} \in \cdot | \Hin_{\to K}=\delta_x),
    \mathbb{P}(\Hin_{\to L} \in \cdot | \Hin_{\to K}=0) \right) \to 0.
  \end{equation}
  where the random variable $\Hin_{\to K}$ is defined as
  \begin{equation}
    \Hin_{\to K}(\omega)= \sum_{i\ge 1} \delta_{X_{\tau_K(w_i)}}
  \end{equation}
  if $\omega = \sum_{i\ge 1} \delta_{w_i^*}$ and  $w_i$ is a representative of $w_i^*$ for each $i\ge 1$. 
  Now,
  \begin{itemize} 
    \item  $\mathbb{P}(\Hin_{\to L} \in \cdot | \Hin_{\to K}=0)$ is the law of Poisson point process on $\mathbb{V}$ with intensity measure $e_L^K (x') := e_L(x')P_{x'}[\tau_K=\infty]$,
    \item while $\mathbb{P}(\Hin_{\to L} \in \cdot | \Hin_{\to K}=\delta_x)$ is the law of the superposition of this Poisson point process and an independent point in $\mathbb{V}$ sampled under $p_L^{K, x} (x') := P_x[X_{\lambda_L}=x'|\tau_K^+=\infty]$.
  \end{itemize}
  By Lemma \ref{lem:general_cond_dTVto0_for_PPP}, it is enough to establish that for every $\epsilon>0$, 
\begin{equation}\label{eq:sum_pLKxy-epshLK_to0_weak}
  \sum_{x'  \in  L} 
  \left(p_L^{K, x}(x') 
  - \epsilon \, e_L^{K}(x') \right)_+ \underset{L \to \mathbb{V}}{\longrightarrow}.
\end{equation}
By reversibility:
\begin{equation}
  a_x \p_x[X_{\lambda_L}=x', \tau_K^+=\infty]  = a_{x'} \p_{x'}[\tau_L^+=\infty] 
   \p_{x'}[\tau_K < \infty, X_{\tau_K} = x]
\end{equation}
hence
\begin{equation}
  p_L^{K, x}(x')= \frac {e_L(x')}{ e_K(x)} 
  \p_{x'}[\tau_K < \infty, X_{\tau_K} = x],
\end{equation}
and therefore, playing on $\epsilon$, \eqref{eq:sum_pLKxy-epshLK_to0_weak} is equivalent to the fact that for every $\epsilon>0$, 
\begin{equation}
    \sum_{x'\in   L}  e_L(x')  \big(  P_{x'}[\tau_K<\infty, X_{\tau_K}=x]  -\epsilon P_{x'}(\tau_K=\infty) \big)_+  \underset{L \to \mathbb{V}}{\longrightarrow}.
\end{equation}
Bounding $P_{x'}[\tau_K<\infty, X_{\tau_K}=x]$ by $P_{x'}[\tau_K<\infty]$, the latter follows from the fact that for every $\epsilon>0$,
\begin{equation}
    \sum_{x'\in  L} e_L(x') \left((1+\epsilon) P_{x'}[\tau_K<\infty]-\epsilon \right)_+  \underset{L \to V}\longrightarrow 0,
  \end{equation}
which, finally, can be derived elementary from Lemma \ref{lem:quantitative_weak} by using the bound:
\begin{equation}
  P_{x'}[\tau_K<\infty] \le \sum_{z\in K}  P_{x'}[\tau_z<\infty]
\end{equation}
and playing on $\epsilon$.
\end{proof}


\subsection{Proof of the 0-1 law for increasing events}

In this section, we prove Theorem \ref{th:increasing_01law}. We begin with a lemma.

\begin{lemma}\label{lem:for_01_increasing}
  For every finite vertex set $K$, for every admissibe path $\gamma$ inside $K$, for every increasing event $E$ belonging to $\widetilde{\mathcal{T}}_{RI}$, we have:
  \begin{equation}
    \p(E | \In_K=\delta_\gamma) \le \p(E| \In_K=0).
  \end{equation}
\end{lemma}


\begin{proof}[Proof of Lemma \ref{lem:for_01_increasing}]
We consider a finite vertex set $K$, and an admissibe path $\gamma$ inside $K$, say with extremities $(x,y)\in \partial_X K\times \partial_X K$. Let furthermore $\epsilon >0$. We are going to construct a vertex set $L$ (containing $K$) and, under a probability denoted $\mathbf{P}$, a coupling $(\omega , \omega')$ of $\p( \cdot | \In_K=\delta_\gamma) \le \p( \cdot | \In_K=0)$ together with random trajectories $w^*, w_1^*, w_2^*\in W^*$ such that, with probability not smaller than $1-\epsilon$: 
\begin{itemize}
  \item $w^*$ belongs to $\omega$ and is precisely the trajectory of $\omega$ that visits $K$, $w_1^*$ and $w_2^*$ belong to $\omega'$;
  \item $\omega- \delta_{w^*} = \omega'-\delta_{w_1^*}-\delta_{w_2^*}$;
  \item $Out_{\to L}(w)=Out_{\to L}(w_1) $, i.e., $w^*$ and $w_1^*$ coincide up to their entry in $L$, in particular their entry points in $L$ coincide;
  \item $Out_{L \to}(w)=Out_{L \to}(w_2) $, i.e., $w^*$ and $w_2^*$ coincide from their last visit to $L$, in particular their last visit points in $L$ coincide;
\end{itemize}
where $w, w_1$ and $w_2$ are representatives of $w^*$, $w_1^*$ and $w_2^*$ respectively. 

\smallskip

To construct the coupling, we use twice the result stated in \eqref{eq:criterion_d_TV_cond_bis} and established during the proof of Theorem \ref{th:weak_01law}, which we formulate in this way: for a large enough $L$, there is a way, given a Poisson point process $\omega_{entry}$ under $ \p( \Hin_{\to L} \in \cdot | \In_K=0)$ to pick a point $X$ in $\omega_{entry}$ such that the couple $(\omega_{entry}-\delta_X, X)$ is at total variation distance at most $\epsilon/2$ from the product law $ \p( \Hin_{\to L} \in \cdot | \In_K=0) \otimes p_L^{K,x}$, and similarly with $x$ replaced by $y$.

So, consider the Poisson point process $\omega'$ distributed under $\p( \cdot | \In_K=0)$. Among the entry points to $L$ of the trajectories in $\omega'$, we pick a point $X$ such as prescribed above, and denote the corresponding trajectory $w_1^*$. Among the exit points of $\omega' - \delta_{w_1^*}$, also pick a point $Y$ such as described above and denote the corresponding trajectory $w_2^*$. Furthermore, sample (under the law of the MC $X$) a trajectory inside $L$ that links $X$ and $Y$ and leaves trace $\eta$ on $K$. Set $w^*$ to be the trajectory that coincides with $w_1^*$ upp to its entry in $L$, is given inside $L$ by the previously sampled trajectory, and coincides with $w_2^*$ after its last visit to $L$. Finally set $\omega$ to be $\omega'-\delta_{w_1^*} - \delta_{w_2^*} + \delta_{w^*}$. 

%choisir autre notation que $X$ et $Y$ ?


% Now, attach to this process $\omega_{entry}-\delta_X$ infinite trajectories according to the random interalcement rules, and consider the exit points of $L$, these form a point process $\omega_{exit}$ which is distributed under $ \p( \Hin_{\to L} \in \cdot | \In_K=0)$ and similarly, there is a way to pick a point $Y$ in $\omega_{exit}$ such that the couple $(\omega_{exit}-\delta_Y, Y)$ is at total variation distance at most $\epsilon/2$ from the law $ \p( \Hin_{\to L} \in \cdot | \In_K=0) \otimes p_L^{K,y}$.
% Finally sample a piece of trajectory from $x$ to $X$ and a piece of trajectory from $y$ to $Y$. This yileds a trajectory $w$ going through $K$. 
% Set $\omega$ to be the process with all the infinite trajectories. 

%Let us assume for now that this coupling exists and derive the lemma. 

\medskip

Let us finally derive the lemma. For every realisation of the coupling, we pick (arbitrarily) a trajectory $w_{add} \in W_L$ such that $Out_{\to L}(w_{add})=Out_{\to L}(w_2) $ and $Out_{L \to}(w_{add})=Out_{L \to}(w_1) $ and we set $w_{add}^* =\pi(w_{add})$. 
We observe that, when the coupling is successful, we have $\widetilde{\Out}_L(\omega +\delta_{w_{add}^*})= \widetilde{\Out}_L(\omega')$. It follows that for every increasing event $E$ belonging to $\widetilde{\mathcal{T}}_{RI}$,
\begin{equation}
  \begin{aligned}
    \p(E | \In_K=\delta_\gamma) 
    & = \mathbf{P}(\omega  \in E)\\
    & \le \mathbf{P}(\omega +\delta_{w_{add}^*} \in E)\\
    & \le \mathbf{P}(\omega' \in E) + \epsilon\\
    & = \p(E| \In_K=0) +\epsilon.
  \end{aligned}
\end{equation}
Since $\epsilon$ is arbitrary, the lemma follows.
\end{proof}


\begin{proof}[Proof of Theorem \ref{th:increasing_01law}]
  From Lemma \ref{lem:for_01_increasing}, we derive that for every finite vertex set $K$, for every admissible configuration $\eta$ inside $K$, for every admissibe path $\gamma$ inside $K$, for every increasing event $E$ belonging to $\widetilde{\mathcal{T}}_{RI}$, we have:
  \begin{equation}\label{eq:extend_for_01_increasing}
    \p(E | \In_K=\eta + \delta_\gamma) \le \p(E| \In_K=\eta).
  \end{equation}
  This relies on the observation that if $E$ is increasing and belongs to  $\widetilde{\mathcal{T}}_{RI}$, then for every (finite) configuration $\sum_{i=1}^k \delta_{w_i^*}$, with $k\in \N$, and $w_i^*\in W^*$, the event
  \begin{equation}
    \left\{ \omega \in \Omega : \omega + \sum_{i=1}^k \delta_{w_i^*} \in E \right\}
  \end{equation} 
  is also increasing and also belongs to $\widetilde{\mathcal{T}}_{RI}$. Just apply this fact to the configuration arising from $\eta$ to derive \eqref{eq:extend_for_01_increasing} from Lemma \ref{lem:for_01_increasing}.

  By induction, it follows from \eqref{eq:extend_for_01_increasing} that for every finite vertex set $K$, for every admissible configuration $\eta$ inside $K$, for every increasing event $E$ belonging to $\widetilde{\mathcal{T}}_{RI}$, we have:
  \begin{equation}\label{eq:extend_encore_for_01_increasing}
    \p(E | \In_K=\eta ) \le \p(E| \In_K= 0 ).
  \end{equation}
  We finally observe that the converse inequality holds, since $\p( \cdot | \In_K=\eta ) $ dominates $\p(\cdot | \In_K= 0 )$. So we have the equality in \eqref{eq:extend_encore_for_01_increasing} and we deduce that any increasing event $E$ belonging to $\widetilde{\mathcal{T}}_{RI}$ is independent of $\In_K$ for every finite vertex set $K$, and therefore is independent of $\mathcal{A}$, thus also of itself, so that $\p(E)\in \{0, 1\}$. 
\end{proof}






\bibliographystyle{plainnat}
\bibliography{bibliography}

\end{document}


********************************************************
********************************************************
********************************************************
********************************************************
********************************************************
********************************************************
********************************************************


A checker : 

AVEC LE NUMERO 888 %888
J'AI INDIQUe dans l'article des modifs à effectuer pour la V2.


%remplacer l'espace $V$ par $\mathbb{V}$ 

%Ecrire FKG et non harris-FKG.

%Cooriger $\mathcal{W}$ doit devenir $\mathcal{W}^+$ à cettains endroits.



DANS LA V2:

- supprimer certains numéros déquations inutiles
- numéroter autrement les théorèmes, elmmes, etc : par section ?

- mettre des illustrations des preuves et des définitions !
- une section sur les conditions équivalentes de tail-triviality et atomicity, notamment en termes de frontière de Martin  - DANS LA V2

- écrire plus joliment 0-1.

- $P[...]$ et non  $P(...)$.



- écrire une remarque qu'on peut coupler les entrelacs avec des configurations différentes prescrites dans K, si la loi du 0 - 1 est valable.


- Une section d'examples

- exemple simple ou loi du 0-1 non vérifiée 
%- l'énoncé et la preuve de FKG, avec peut-être une preuve générique (ou schéma) pour PPPP satisfait FKG, et lien avec la preuve de Teixeira.

- applications des lois de 0-1. Par exemple pour graphe fois $\Z$: critère quantitatif fonctionne toujours ? adaptation critère de couplage.

- application de FKG (ou 0-1)  pour définir paramètre critique.


- Remarque: difficulté à généraliser la pruev deu cas atomqiue au cas non atomique, choses à comprendre, notamenet du poitnt de vue de la théorie de la frontière de Martin. cf. idéées dans le document Lyx "Notes et idées de recherche".


- réfléchir à ce qu'il se passe si le graphe n'est pas localement fini.

- réfléchir aux entrelacs si $X$ n'est pas réversible. Les parties de $\omega$ qu'on obtient en restreignant à $A \to B$ sont elles de ce type ?



\section{Examples and applications}

\subsection{Examples}

\begin{example}\label{ex:simple_not_01_copy}%IN V2: send to section examples.
Consider the Markov associated with the nearest neighbor graph on $\Z$ and the family of conductances:
\begin{equation}
  a_{n, n+1}=2^n, \qquad a_{-n-1, -n}=2^n, \qquad  n\ge 0, 
\end{equation}
i.e., the Markov chain with the transition probabilities $p_{0, -1}=p_{0, 1}=\frac{1}{2}$, and, for $n\ge 1$, $p_{n, n+1}=p_{-n, -n-1}=\frac{2}{3}$ and $p_{n, n-1}=p_{-n, -n+1}=\frac{1}{3}$.
Let us consider the trajectories of $\omega$ that go from $-\infty$ to $\infty$. These trajectories have to pass through 0, so their number is distributed as a Poisson variable with parameter 
\begin{equation}
  \begin{aligned}
  \lambda:=& P_0[X_n \underset{n\to\infty}\longrightarrow -\infty |\tau_0^+=\infty ] e_{\{0\}}(0) P_0[X_n \underset{n\to\infty}\longrightarrow \infty]\\
  = & P_0[X_n \underset{n\to\infty}\longrightarrow -\infty, \tau_0^+=\infty ] a_0 P_0[X_n \underset{n\to\infty}\longrightarrow \infty ].
  \end{aligned}
\end{equation}
The parameter $\lambda$ is clearly positive and finite% (it is in fact equal to $1/4$)
, so the number of trajectories in $\omega$ going from $-\infty$ to $\infty$ is a nontrivial random variable. It is however measureable with respect to $\mathcal{T}_{RI}$. 
\end{example}

in fact one could just ask the conductances to be symmetric (and the MC to be transient), so that the two escapes are clearly possible.


More genral example with simply the condition that $\sum_{n\ge 0} \frac{1}{a_{n, n+1}}=\infty$ and $\sum_{n\le 0} \frac{1}{a_{n, n+1}}=\infty$ so that both escapes are possible.

Example of the ternary (or $d$-regular) tree with conductances 1: no 0-1 law. Distance to a fixed vertex evolves as the first example.

More generally example of trees: probas can be expressed. Recall formulas.


Produit avec Z d'un graphe infini. On met des arêtes $(x,n)\sim (x, n+1)$ avec capacité $a_x/2$. Cela produit la chaîne produt de $G$ et $\Z$. 
Ou bien les arêtes : $(x,n)\sim(y,m)$ avec proba $p_{x,y} p_{n, m}$: produit indépendant des deux chaînes. Maiière très artificielle de produire beaucoup de trajectoires.
Question: est-ce que la loi 0-1 est toujours vérifiée ?


Graphe sur $\Z$ (de l'exmaple 1, par exemple) multipplié par $\Z$: toujours deux issues mais les paramètres $Q(A\to B)$ sont tous infinis.

Arbre ternaire produit $\Z$: Hors du cas atomique. critère quantitatif vérifié ? Condition nécessaire $Q(A \to B) $ vérifiée ? 
Possibilité d'adpter la preuve de cas atomique?

Question : quels liens avec loi 0-1 pour événements invaraints par translation. Notre approche plus robuste a priori : on peut faire de smodifs locales sans affecter les critères ?? Pas si clair... 

Remark : for the trinary tree, a 0-1 law holds, with respect to translation invariant events. 






\subsection{Applications}

Peu sensible aux modifications locales du graphe pondéré ?


Pour I, on a presque sûrement $|I|$ est infini, p.s., toutes les composantes connexes de $I$ sont infinies.


\subsection{APPLICATIONS LOI DU 0-1 pour événements croissants}

de plus, l'événement "I possède une infinité de composantes connexes" est de proba 0 ou 1.

Pour V, les événements suivants ont proba 0 ou 1:

$|V|=\infty$, $V$ a une infinité de composantes connexes. V a au moins une composante connexe infinie. V a une infinité de composantes connexes infinies. 

A propos de $\omega$:
 les événements suivants sont de proba 0 ou 1:
 il existe une infiité de sites $x$ avec temps local en $x$  au moins 10.

 il existe une infiité de $x$ tel que $B(x,10)\subset V$. idem pour I.

 Remplacer 10 par n'importe quelle famille $(r_x)_x$.


 

Remarque : si l'on a l'existence d'un $\epsilon>0$ et une collection $(B_i)_{i\in I}$ d'ensembles dans $V$ telle que $B_i$ s'loigne de $0$ et $P[B_i \cap V \neq \emptyset]\ge \epsilon$ alors p.s. $V$ est infini.
C'est l'argument utilisé avec Serguei pour le 2DRI.

\subsection{application loi du 0-1 (si elle est valable)} 

 les événements suivants sont de proba 0 ou 1:

Il existe une infiité de $x$ tel que le temps local en $x$ est entre $10$ et $20$.


 il existe une infiité de $x$ tel que $B(x,10)\subset V$ mais pas $B(x, 20)\subset V$. idem pour I.

 
 Remplacer 10 et 20 par n'importe quelles familles $(r_x)_x$ et $(R_x)$.

etc... 



Par contre; le nombre de composantes connexes ou le nombre de composantes connexes infinies ne sont pas non-locaux donc pas de résulatt de ce côté là.


Il faut chercher des idées d'appliations, par exemple pour des modèles particuliers.


Idée d'application: vacant set dans un "random conductance model" sur $Z^d , d\ge 3$. Déjà prouvable par invaraince par translation de la loi du graphe ? Qu'en est-il si on fait des modifs locales dans la loi.

\appendix

\section{About $\mathcal{T}_{MC}$}



We now provide if and only if conditions for the occurence of the 0-1 law for $\mathcal{T}_{MC}$ (defined in ???).

\begin{lemma}\label{lem:equiv_conditions_01_MC_appendix}
The following conditions are equivalent:
\begin{itemize} 
  \item[(i)] The Markov chain $X$ satisfies the 0-1 law for $\mathcal{T}_{MC}$,
  \item[(ii)] For every $x, y\in \mathbb{V}$, $d_{TV}(P_x[X_{\tau_{K^c}} \in \cdot], P_y[X_{\tau_{K^c}} \in \cdot]) \underset{K\to \mathbb{V}}\longrightarrow 0$.
  \item[(iii)] For every $x, y\in \mathbb{V}$, $d_{TV}(P_x[X_{\lambda_K} \in \cdot], P_y[X_{\lambda_K} \in \cdot])  \underset{K\to \mathbb{V}}\longrightarrow 0$.
  \item[(iv)] For every $x, y\in \mathbb{V}$, there exists a coupling $X=(X_n)_{n\ge 0} , Y=(Y_n)_{n\ge 0}$ of the Markov chain $X$ started respectively at $x$ and $y$ such that, almost surely, there exist (random) finite times $S$ and $T$ for which
  \begin{equation}
    \forall n\ge 0 , \qquad  X_{S+n}=Y_{T+n}, 
  \end{equation} 
  that is, $X$ and $Y$ coincide eventually, up to time-shift. 
  \item[(v)] There exists an $\epsilon >0$ such that for every $x, y\in \mathbb{V}$, there exists a coupling $(X, Y)$ (under a probability $\mathbf{P}$) of the Markov chain $X$ started respectively at $x$ and $y$ such that 
  \begin{equation}
    \mathbf{P}[\text{Range}(X) \cap \text{Range}(Y) \neq \emptyset]  \ge \epsilon.  
  \end{equation}
  \item[(vi)] The Markov chain $X$ satisfies the Liouville property: every function $h:\mathbb{V} \to \R$ that is harmonic with respect to $X$ and bounded is constant.
\end{itemize}
\end{lemma}

\begin{proof}
  \begin{itemize}
  \item The implication $(i)\implies(ii)$ follows from Lemma ??. Indeed, if  the 0-1 law for $\mathcal{T}_{MC}$ is satisfied, then the Markov chain $(X_{\tau_{K_n^c}})_{n\ge 0}$, where $(K_n)_{n\ge 0}$ is a fixed exhaustion of $\mathbb{V}$, is tail-trivial and an application of Lemma ?? yields $(ii)$.
  \item The implication $(ii)\implies (iii)$ follows from the bound
  \begin{equation}
    d_{TV}(P_x[X_{\lambda_K} \in \cdot], P_y[X_{\lambda_K} \in \cdot])
    \le 
    d_{TV}(P_x[X_{\tau_{L^c}} \in \cdot], P_y[X_{\tau_{L^c}} \in \cdot])  
  \end{equation}
  which holds as soon as $K$ contains the set $\{x: \exists y\in L,  x\sim y\}$.
  \item The implication $(iii)\implies (iv)$ relies on a maximal coupling. Assume that $(iii)$ holds and pick $x, y\in \mathbb{V}$. Consider an exhaustion $(K_n)_{n\ge 0}$ of $\mathbb{V}$ such that $K_0$ contains $x$ and $y$. The process $(Z_n)_{n\ge 0}:=(X_{\lambda_{K_n}})_{n\ge 0}$ is a Markov chain. Consider a maximal coupling of this process, started from $P_x(X_{\lambda_{K_0}}\in \cdot )$ and $P_y(X_{\lambda_{K_0}}\in \cdot )$, by $(ii)$ the two processes almost surely coincide eventually. Now, complete these to obtain two trajectories of the Markov chain, started respetively from $x$ and $y$. Where the two copies of $Z$ coincide, build the completion identically in bot copies. Then, almost surely, the two trajectories eventually coincide up to time-shift.
  \item The implication $(iv)\implies (v)$ is trivial.
  \item We show the implication $(v)\implies (i)$ by contraposition. So we assume that $(i)$ does not hold and show that $(v)$ is not satisfied. We pick $E$ and $F$ two disjoint events in $\mathcal{T}_{MC}$ having positive probability under $P_x$ (for arbitrary $x$). For every $p \in [\frac{1}{2}, 1)$, let us denote 
  \begin{equation}
    E_p :=\{ x \in \mathbb{V} : P_x(E) > p\}\, , \qquad F_p :=\{ x \in \mathbb{V} : P_x(F) > p\}
  \end{equation}
  Since the events $E$ and $F$ are disjoint, we see that $E_{1/2}$ and $F_{1/2}$ are disjoint. 
  Let us show that for every $p\in [\frac{1}{2}, 1)$, $E_p$ and $F_p$ are non-empty.  COMPLETE

  Now, for every $p\in (1/2, 1)$, for every $x\in E_p$, we observe, using the Markov property at time $\tau_{(E_{1/2})^c}$ and the fact that $E\in \mathcal{T}_{MC}$, that
  \begin{equation}
  1- p > P_x (E^c) \ge P_x (\tau_{(E_{1/2})^c} <\infty , E^c) \ge P_x (\tau_{E_{1/2}^c} <\infty ) \times \frac{1}{2},
  \end{equation} 
  so that $  P_x(\tau_{(E_{1/2})^c} <\infty ) < 2(1-p)$.   A similar result holds for $F$. 

  Let us finally prove that $(v)$ does not hold. Let $\epsilon \in (0, 1)$. Let us pick an $x\in E_{1-\epsilon/4}$ and an $y\in E_{1-\epsilon/4}$. Recall that $E_{1/2}$ and $F_{1/2}$ are disjoint. 
  For any coupling  $(X,Y)$ of the Markov chain $X$ started respectively from $x$ and $y$, if the ranges of $X$ and $Y$ intersect, then $X$ must leave $E_{1/2}$ or $Y$ must leave $F_{1/2}$, but the probabilities of these two events are smaller than $\epsilon/2$, so the ranges of $X$ and $Y$ intersect with probability smaller than $\epsilon$. Thus $(v)$ does not hold.
  \item We show the implication $(i)\implies (vi)$. Let us assume that $(i)$ holds and let $h:\mathcal{V}\to\R$ be a bounded and harmonic (with respect to $X$) function. 
  Consider $x\in \mathbb{V}$. The process $(h(X_n))_{n\ge 0}$ is a bounded martingale, so, under $P_x$, $h(X_n)$ converges towards a random variable. This limiting random variable is measurable with respect to $\mathcal{X}$ so it must be almost surely constant. By the stopping theorem, $h(x)=E_x[h(X_n)]$ for every $n\ge 0$, hence the limiting random variable is equal to $h(x)$. We have established:
  \begin{equation}
    \forall x\in \mathbb{V}, \qquad P_x-\text{a.s.}, \quad h(X_n ) \underset{n\to\infty}\longrightarrow h(x). 
  \end{equation}
  Consider now $x$ and $y$ in $\mathbb{V}$. Let us consider, under $P_x$, the event that $X$ visits $y$ (which has positive probability by irreducibility of $X$). On this event, we have that $h(X_n)$ must converge toward $h(x)$ but also, applying the Markov property at time $\tau_y$, that $h(X_n)$ must converge toward $h(y)$. It follows that $h(x)=h(y)$. We have shown that $h$ is constant. 
  \item We show the implication $(vi)\implies (i)$. Assume $(vi)$ and consider $E$ an event in $\mathcal{T}_{MC}$. The function $h:\mathbb{V}\to\R, x\mapsto P_x(E)$ is bounded and harmonic with respect to $X$, so it must be constant, let us denote its value by $p_0\in[0, 1]$. Let us assume that $p_0>0$. Then, we may use the result show in ??: for every $p\in [1/2, 1)$ the set $\{x\in \mathbb{V} : P_x(E) >p\}$ is non empty. It follows that $p_0=1$. 
\end{itemize}

  (iii) implies (ii) is  clear. (ii) implies (iii) follows from maximal coupling. 

  Let us prove that (ii) implies (i). Let $A$ be such a tail event. Consider any finite path $(x_0, \dots, x_l)$ having positive probability under $P_x$, i.e., $P_x[X_0=x_0, \dots, X_l=x_l]>0$. 
  We are going to show that, under $P_x$, the tail event $A$ is independent of the event $\{X_0=x_0, \dots, X_l=x_l\}$. Let $\epsilon >0$. Consider $K$ large enough so that it contains the path $(x_0, \dots, x_l)$ and so that 
  $d_{TV}(P_x[X_{\tau_{K^c}} \in \cdot], P_{x_l}[Y_{\tau_{K^c}} \in \cdot]) \le \epsilon$. 
  Using the Markov proerty, the Markov chain started from $x$ and conditioned on the event $\{X_0=x_0, \dots, X_l=x_l\}$ is distributed after time $l$, as the Markov chain started from $l$. We can thus couple it with the Markov chain started from $x$ such that there entry point to $K^c$ coincide with probability not smaller than $1-\epsilon$. Since $A$ is determined by the behaviour of the walk outside $K$, we deduce:
  \begin{equation}
    |P_x(A)-P_x(A|X_0=x_0, \dots, X_l=x_l) | \le \epsilon.
  \end{equation}
  Since $\epsilon$ is arbitrary, we obtain the desired result: $P_x(A|X_0=x_0, \dots, X_l=x_l)=P_x(A)$. We derive that $A$ is independent of any cylindrical event, hence is independent of the whole $\sigma$-algebra W [write correctly], in particular it is independent of itself and therefore as probability either 0 or 1.
  
  Let us finally prove that (i) implies (ii). We reason by contradiction. Let us assume that the 0-1 law holds and that there exist $x$ and $y$ such that
  \[\lim_{K\to V} d_{TV}(P_x[X_{\tau_{K^c}} \in \cdot], P_y[Y_{\tau_{K^c}} \in \cdot]) :=\delta >0 \]
  (the limit exists by monotonicity, which follows from elementary arguments). We fix $(K_n)$ an exhaustion of $V$. Let us consider the maximal coupling started from $x$ and $y$ and for each $n$ the subset $A_n$ of $V$ defined as:
  \[A_n:=\{z \in V, P_x[X_{\tau_{K_n^c}} = z] > P_y[X_{\tau_{K_n^c}} = z] \} \,, \]
  which is included in $K_n^c$ and satisfies 
  \[P_x[X_{\tau_{K_n^c}} \in A_n] > P_y[X_{\tau_{K_n^c}} \in A_n] +\delta \ge \delta. \] 
  by definition of the total variation distance and by  [eq. def of delta].
  Let us consider the event $A=\{X_{\tau_{K_n^c}}\in A_n for infinitely many n's\}$. It is a tail event, and thus has probability belonging to $\{0, 1\}$, for the Markov chain started from $x$ as well as for the Markov chain started from $y$: in fact the value is the same in both cases, let us denote it by $p$. If $p=0$, i.e., $P_x[X_{\tau_{K_n^c}}\in A_n for infinitely many n's]=0$, then $P_x[X_{\tau_{K_n^c}}\in A_n   ]$ goes to 0 as $n\to\infty$, which contradicts [eq.]. So $p$ must be 1 and therefore almost surely, the Markov chain started from $y$ will satisfy $X_{\tau_{K_n^c}} \in A_n$ for some $n$ (in fact, infinitely many), but by maximality of the coupling, this implies that the Markov chains have already coupled and they will remain together afterwards. Then (iii) is established, and therefore (ii) also. By contradiction (i) implies (ii). 
\end{proof}

??? OBSERVE VALID ALSO WHEN  NOT REVERSIBLE?





\begin{remark}\label{rem:tail_atomicity_Martin}???
tail atomicity condition oin terms on Martin Boundary. SEND THIS TO APPENDIX.
\end{remark}



\bibliographystyle{plainnat}
\bibliography{bibliography}

\end{document}


**************************************************************
**************************************************************
**************************************************************
**************************************************************


% We finish this section by showing the identity mentioned in Remark ??, that is: when $x,y\in L$,

% \begin{equation}
%   P_{x'}[\tau_x < \infty | X_{\lambda_L}=y']= P_{x'}[\tau_x < \tau_{y'} | \tau_{y'}<\infty].
% \end{equation}

% We write, using several times the strong Markov property:
% \begin{equation}
%   \begin{aligned}
%   P_{x'}[\tau_x < \infty | X_{\lambda_L}=y']  
%   & = \frac{ P_{x'}[\tau_x < \infty, X_{\lambda_L}=y'] }{P_{x'}[X_{\lambda_L}=y']} \\
%    & = \frac{ P_{x'}[\tau_x < \infty] P_x[X_{\lambda_L}=y']}{P_{x'}[X_{\lambda_L}=y'] } \\
%    & = \frac{ P_{x'}[\tau_x < \infty] P_x[\tau_{y'}<\infty ] P_{y'} [X_{\lambda_L}=y'] }{P_{x'}[\tau_{y'}< \infty] P_{y'} [X_{\lambda_L}=y'] } \\
%    & = \frac{ P_{x'}[\tau_x < \infty] P_x[\tau_{y'}< \infty] }{P_{x'}[\tau_{y'}< \infty] }
%   \end{aligned}
% \end{equation}  




**************************************************************


% (((
  
% We introduce tail $\sigma$-algebras, for the Markov chain and for the random interlacements process. For the Markov chain $X$, we do not use the usual tail $\sigma$-algebra, we rather consider a notion of tail $\sigma$-algebra that is adapted to the equivalence up to time shift.
% HOW TO DEFINE IT? PRObably MULTIPLE WAYS... 
% Note that is contained in the usual sigma-algebra. 

% Def 1: events in the usual that are invaraint under time shift (n,ote that they belong to the usual tail sigma algebra)

% Def 2: intersection of the sigma algebras generated by the pieces of trajectory oustide K, K going to V. (oustide meaning after the first vsit to complementray or last visit to K). 

% Def 3 ?: Quotient space under equivalence up to time shift ?

% DEF 1 and 3: just look at events that are invariant under time shifts. Are they automatically tail events in the classical sense?

% )))

% First, the tail $\sigma$-algebra for the Markov chain $(X_n)_{n\ge 0}$ is the sub-algebra of ?? defined as usual as:
% \begin{equation}
%   \mathcal{T}_{MC} = \bigcap_{n\ge 0} \sigma(X_m,  m\geq n).
% \end{equation}




**************************************************************



  
%   We consider $x\in V$ and $\epsilon>0$. 
%   We consider the set  
% $D=\{y \in V : P_y[\tau_x< \infty] > \epsilon\}$. Let us consider first a very convenient situation, which is when $D$ is finite. Then, also $\bar{D}=\{y\in V : d(y, D)\le 1\}$ is finite and for any $L$ containing $\bar{D}$, the sum in Lemma ?? is zero: indeed, any $y$ in the interior frontier of $L$ is outside $D$, so satisfies $P_y[\tau_x< \infty]\le \epsilon$.

% It is however not true in general that $D$ is finite (going against the naive idea that the probability of hitting $x$ should decay to 0 as the starting point goes away from $x$). See counterexamples in ??. Since the criterion in lemma ? is restricted o finite vertex sets $L$, the argument needs to be "repaired" (find better formulation: patched ?).


% To handle this situation we observe that $D$ is transient. Indeed, a trajectory which visits $D$ infinitely often would also, by simple arguments, almost surely hit $x$ infinitely often, which contradicts the transience of the chain. 
% We also observe that the total capacity of $D$ (defined as the mass $\nu(W_D)$ is finite: indeed, $\capa(\{x\}) \ge \epsilon  \capa(D)$ (check it through an exhaustion of $D$).

% So a bi-infinite trajectory of the interlacement process conditioned on hitting $x$ (so under $\nu$ restricted to $\{x\}$) almost surely has an entrance and an exit point to $D$. 

% Let us consider $(L_n)$ an exhaustion of $V$ (and assume that $x$ belongs to $V_1$ for simplicity). 
% We can separate the frontier of $L_n$ into two subsets:
% \begin{itemize}
% \item The sites $x'$ belonging to $V\setminus D$: there, $\left(P_{x'}[\tau_x < \infty]-\epsilon \right)_+ =0$ by defintion of $D$
% \item The sites $x'$ belonging to $D$: there, we bound $ \left(P_{x'}[\tau_x < \infty]-\epsilon \right)_+ $  by 1.
% \end{itemize}
% Hence we bound the sum in ?? by $e_{L_n}((\partial_{in} L_n) \cap D)$.

% This is the total mass, under $\nu$, of the trajectories that visit $L_n$, and enter it through $D$. However, 

% But the total mass of $e_{L_n}$ converges to $\capa(D)$ and $e_{L_n}((\partial_{in} L_n) \cap (V\setminus D))$ converges also towards $\capa(D)$ since the rajectories almost surely enter $D$ at some point. UNCLEAR, REPHRASE.


% For a bi-infinite trajectory of the interlacement process conditioned on hitting $x$ (so under $\nu$ restricted to $\{x\}$ renormalised), almost surely, for large enough $n$, it enters $L_n$ through $ \partial L_n \cap(V\setminus D)$. 


% It is maybe best to consider the explicit exhaustion :

% $L_n:= \text{ blow up by 1 of } \{y: d(x, y)\le n, y\in D\}$.

% This is an exhaustion of $\bar{D}$, not $V$ so it must be also blown in some way... 



**************************************************************
**************************************************************
**************************************************************


\part{Vieux brouillon}


\section{Interpretation of RI as a Poisson loop soup}

Section pas forcément nécessaire.

Assumption is that for $x$ and $y$ away from some reference point 0, the MC started from $x$ and from$y$ conditioned to touch $K$ can be coupled so that their entry points to $K$ coinicde with large probability. and then the law is $harm_K$.


Has some links with the assumption that MC started from $x$ and $y\in K$ and conditioned to escape $K$ can be coupled to coincide eventually.








********************************************************
********************************************************
********************************************************
********************************************************

Let us denote $\N=\{1, 2, \dots\}$, $\N_0=\N\cup\{0\}$.


\section{Poisson point processes satisify the FKG property}

Let $(X, \mathcal{X})$ be a measurable space. We define the following set of measures on $(X, \mathcal{X})$
\[\Omega : = \{ \omega = \sum_{n\in \N} \delta_{x_n}, (x_n)_{n\in \N} \in X^\N \}\]
and equip it with the $\sigma$-algebra $\mathcal{A}$ generated by the applications $\omega \mapsto \omega(D)$, for $D\in \mathcal{X}$.
We define an order on $\Omega$:
\[\omega \le \omega' \quad \text{ if and only if } \quad \forall D \in \mathcal{X},  \omega(D) \le \omega'(D).\]

Let $\mu$ be a $\sigma$-finite measure on $(X, \mathcal{X})$ 
(CONDITIONS ON $\mu$? : YES Radon measure, so the space must be equipped with a topology).
There exists a law on $(\Omega,\mathcal{A})$, denoted $PPP(\mu)$ and called law of the (?) Poisson point process with intensity measure $\mu$, such that under $PPP(\mu)$:
\begin{itemize}
\item for any $D\in \mathcal{X}$, $\omega(D)$ is a Poisson variable with parameter $\mu(D)$,
\item for any $n\in\N$, any disjoint $D_1, \dots, D_n \in \mathcal{X}$, the variables $\omega(D_1), \dots, \omega(D_n)$ are independent. 
\end{itemize}

We may sometimes identify $\Omega$ with ${\N_0}^X$, by identifying $\omega$ with the family $(\omega(\{x\}))_{x\in X}\in {\N_0}^X$. The order on $\Omega$ then corresponds to the product order on ${\N_0}^X$. 

In the case where $X$ is countable, we will always take $\mathcal{X}$ to be the discrete $\sigma$-algebra on $X$; in that case, the $\sigma$-algebra $\mathcal{A}$ corresponds to the product $\sigma$-algebra on ${\N_0}^X$, and $PPP(\mu)$ corresponds to the product law $\bigotimes_{x\in X} Poi(\mu(x))$.


\begin{prop}
We assume that $X$ is countable and we consider a family $(\mu_x)_{x\in X}$ of measures on $(\N_0, \mathcal{P}(\N_0)) $. 
Then the product law $\Pi:=\bigotimes_{x\in X} \mu_x$, defined on $\N_0^X$ equiped with the product $\sigma$-algebra satisfies the FKG-inequality: 
for any measurable non-decreasing functions $f, g : \N_0^X \to \R$, we have
\[\E^{\Pi}[f g ] \ge \E^{\Pi}[f ] \times \E^{\Pi}[g ] . \] 
\end{prop}

\begin{cor}
As a consequence, if $X$ is countable, then $PPP(\mu)= \bigotimes_{x\in X} Poi(\mu(x))$ satisfies the FKG-inequality: 
for any measurable non-decreasing functions $f, g : \R^X \to \R$, we have
\[\E^{PPP(\mu)}[f g ] \ge \E^{PPP(\mu)}[f ] \times \E^{PPP(\mu)}[g ] . \] 
\end{cor}

We are now going to remove the hypothesis that $X$ is countable.


\begin{prop}
The law $PPP(\mu)$ satisfies the FKG-inequality: for any measurable non-decreasing functions $f, g : \Omega \to \R$, we have
\[\E^{PPP(\mu)}[f g ] \ge \E^{PPP(\mu)}[f ] \times \E^{PPP(\mu)}[g ] . \] 
\end{prop}

\begin{proof}


If $Z$ is a subset of $\mathcal{X}$, we denote by  $\mathcal{F}_(Z)$ the $\sigma$-algebra  generated by the applications $\omega  \mapsto \omega(D), D\in Z$.
We say that $A\in \mathcal{A}$ is finite-dependent if it belongs to $\mathcal{F}_Z$, for some finite subset $Z$ of $\mathcal{X}$. 

To begin with, we handle the case when $f$ and $g$ are finite-dependent, say they are measurable with respect to $\mathcal{F}_Z$, with $Z=\{D_1, \dots, D_n\}$ is a finite subset of $\mathcal{X}$. 
In fact, without loss of generality, we can assume that the $D_i$'s are disjoint: otherwise, it suffices to replace them by all the sets of the form $\cap_{i=1}^n D_i^{c_i}$, for $(c_i)_{i=1, \dots, n} \in \{1, 2\}^n$, and denoting $D_i^1 = D_i, D_i^2= X\setminus D_i$. 

It follows, by definition, that under $PPP(\mu)$, the variables $\omega(D_1), \dots, \omega(D_n)$ are independent. 
Since $f$ and $g$ are non-decreasing functions of these variables, applying Proposition (?), we get the desired result:
\[\E^{PPP(\mu)}[f g ] \ge \E^{PPP(\mu)}[f ] \times \E^{PPP(\mu)}[g ] . \] 


\medskip 

For the rest of the proof, we use classical reasonings in measure theory. The aim is to extend (??) to all non-decreasing measurable functions $f, g$. In fact, we rather aim at showing the following equivalent statement of the FKG inequality: for every increasing events $A$ and $B$, using the short-hand $P=PPP(\mu)$,
\[P(A\cap B)\ge P(A) \times P(B).\]

Let us show the following fact: for any increasing event $A\in \mathcal{A}$, for every $\epsilon>0$, 
there exists an increasing finite-dependent event $A^\epsilon$ contained in $A$ such that $P[A \setminus A^\epsilon] \le \epsilon$.

Let us first show that any event $A\in \mathcal{A}$ can be approximated by finite-dependent events.
To prove this, let us introduce the set 
\begin{align*}
\{A \in \mathcal{A} : \quad & \forall \epsilon>0,  \exists A^-, A^+ \in\mathcal{A}, \text{ finite dependent such that } \\
& A^- \subset A \subset A^+, P(A\setminus A^-)\le \epsilon, P(A^+\setminus A)\le \epsilon\}.
\end{align*}
One can easily check that this set is a $\sigma$-algebra and contains all finite-dependent events. Therefore, it is equal to $\mathcal{A}$.

%For every event $A\in \mathcal{A}$, let us denote the following set:
%\[\text{Under}(A):= \{\omega\in \Omega : \exists \omega'\in A, \omega\le \omega'\},\]
%and let us consider the set
%\[\{A \in \mathcal{A} : \text{Under}(A) \text{ and } \text{Under}(A^\complement) \text{ belong to } \mathcal{A}\}.\]
%One can easily check that this set is a $\sigma$-algebra FALSE! and contains events of the form $\{\omega \in \Omega : \omega(D)\le n\}$, with $D\in \mathcal{X}, n \in \N_0$. Therefore, it is equal to $\mathcal{A}$.

We now prove the claimed property. Let us consider an increasing event $A\in \mathcal{A}$ and $\epsilon>0$. There exists an event $A^-\subset A$ such that $P(A^-)\ge P(A)-\epsilon$. Let us consider the set 
\[\text{Above}(A^-):= \{\omega\in \Omega : \exists \omega'\in A^-, \omega'\le \omega\},\]
In great generality it is not clear that $\text{Above}(A^-)$ is an event, but here $A^-$ is finite-dependent, so it is quite simple to establish that it is the case [DETAILS NEEDED]. Furthermore, it is increasing, contains $A^-$ and is contained in $A$ (since $A$ is increasing). Thus, the point (?) is proved.

\smallskip

We are now ready to conclude the proof. Consider $A$ and $B$ two increasing events. Recall that our aim is to prove that, with $P=PPP(\mu)$:
\[P(A\cap B)\ge P(A) \times P(B).\]
For any $\epsilon>0$, considering $A^\epsilon$ and $B^\epsilon$ as above, we write
\[P(A\cap B)
\ge 
P(A^\epsilon \cap B^\epsilon)
\ge P(A^\epsilon )  \times P( B^\epsilon) 
\ge  (P(A)-\epsilon ) \times (P( B) -\epsilon).\]
Sending $\epsilon$ to 0 yields the desired result.

\end{proof}


\section{Application to Random Interlacements}

Let us consider the point process describing the random interlacement, at some level $u>0$, on a transient weighted graph $G=(V, E,  c)$. It is a Poisson point process on some space  $(W^*, \mathcal{W}^*)$ of trajectories in the graph, which we do not detail here. Hence this Poisson point process satisfies the FKG property. As a consequence, for any ordered set $(\Sigma, \le)$, for any increasing function $\pi : \Omega \to \Sigma$, the variable $\pi(\omega)$ satisfies the FKG property. 

For example, the FKG property is satisfied by the following variables:
\begin{itemize}
\item the collection $\eta=(\eta_x)_{x\in V}\in \{0, 1\}^V$ defined by 
\[\eta_x = 
\begin{cases}
 1 & \text{if } x \text{ is visited by at least one trajectory}, \\
0 & \text{else} 
\end{cases}
\]
(this is the FKG property shown in \cite{Tei09}.)
\item the collection $\eta=(\eta_x)_{x\in V} \in \N_0^V$ defined by 
\[\eta_x = \text{ number of trajectories visiting } x.
\]
\item the collection $\eta=(\eta_x)_{x\in V} \in \N_0^V$ defined by 
\[\eta_x = \text{ number of visits to } x.
\]
\end{itemize}

We can be even more general and consider the full interlacement point process, where trajectories are labelled by positive reals: it is a PPP on $(W^* \times \R, \mathcal{W}^*\otimes \mathcal{B}(\R))$ and it therefore also satisfies the FKG property.





********************************************************
********************************************************
********************************************************
********************************************************
********************************************************
********************************************************


Suppressed chapter of the thesis


For $K$ a finite subset of $\bbS$, we denote 
\begin{equation}
h_K(x) \, = \, P_x[\varphi_K= \infty] \, .
\end{equation}

\begin{lemma}
For every $x\in\bbS$ such that $P_x[ \varphi_K= \infty]>0$, under $P_x[ S \in \cdot | \varphi_K= \infty]$, $h_K(S_n)$ converges in probability towards 1, as $n\to\infty$.

Similarly, for every $x\in\bbS$ such that $P_x[ \forall n \ge 1, S_n \notin K]>0$, under $P_x[ S \in \cdot |  \forall n \ge 1, S_n \notin K]$, $h_K(S_n)$ converges in probability towards 1, as $n\to\infty$.
\end{lemma}

\begin{proof}
Using the Markov property, for every $x, y\in \bbS$, we have
\begin{equation}
P_x[S_n = y , \varphi_K = \infty] = P_x[S_n = y , \varphi_K > n] \times P_y[\varphi_K = \infty] \, ,
\end{equation}
hence
\begin{equation}
\begin{aligned}
E_x\left[\frac{1}{h_K(S_n)} \Big| \varphi_K = \infty \right]
& = \sum_{y \in \bbS} \frac{P_x[S_n = y , \varphi_K = \infty]}{P_x[\varphi_K= \infty]} \times \frac{1}{h_K(y)}\\
& = \sum_{y \in \bbS} \frac{P_x[S_n = y , \varphi_K > n]}{P_x[\varphi_K= \infty]} \\
& = \frac{P_x[\varphi_K >n]}{P_x[\varphi_K= \infty]} \underset{n\to\infty}{\longrightarrow} 1\, .
\end{aligned}
\end{equation}
Using the bounds $0\le 1-x \le \frac{1}{x} - 1$ for $x\in(0, 1]$, we deduce that $h_K(S_n)$ converges towards 1 in $\bL^1$, under $P_x[ S \in \cdot | \varphi_K = \infty ]$. Hence, the convergence holds also in probability.

For the second point in the lemma, use the Markov property at time 1 and then apply the first point.
\end{proof}

\begin{remark}
We observe that lemma ?? is a special instance of theorem ?? in [??].
Extension to almost sure ?
\end{remark}




\begin{figure} 
\begin{center}
\includegraphics[scale=.6]{figures/coupling_RI.pdf}
\caption{Coupling of $\omega$ conditioned on $\In_K$ and $\omega'$ unconditioned. First display : the configuration $\eta$. Second display : $\In_L(\omega)$ (in red, on the left), and $\In_L(\omega)$ (in green, on the right), constructed to have the same number of trajectories (here, 3). Third display : we represent only $\Out_L(\omega)$ and $\Out_L(\omega')$ to keep the figure light : the hinge points of $\omega$ (in red) and of $\omega'$ (in green) come as couples, they are represented with the same symbols. The hinge couples of $\omega$ and $\omega'$ are paired arbitrarily (red crosses with green crosses, etc.) and we couple the trajectories started from the hinge points so that they coincide eventually, up to time shift. After they coincide, they are represented by dashed lines. The trajectories that do not touch $L$ are sampled identically for $\omega$ and $\omega'$; they are represented by dashed lines. ADD A TIME DIRECTION : ARROWS.}
\label{f_coupling}
\end{center}
\end{figure}




\begin{remark}
Let us stress, although we do not need it here, that it is possible to construct an almost sure coupling. 
Indeed, given $K$ and $\eta$, there exists a coupling $(\omega, \omega')$ of $\p( \cdot | \In_K(\omega) = \eta) $ and $\p$ such that, almost surely, $\omega$ and $\omega'$ coincide outside some finite (large, random) subset $M$ of $\bbS$, in the sense that $\Out_M(\omega)= \Out_M(\omega')$. 

To build it, use the same approach as in the above proof. First, couple the processes such that there exists a random $L$ containing $K$ such that the number of trajectories touching $L$ are the same in $\omega$ and in $\omega'$ : this is possible by considering larger and larger $L$'s and using the Borel-Cantelli lemma (see also Remark ?? USING MAXIMAL COUPLING). Once $L$ is obtained, and $\In_L(\omega)$ and $\In_L(\omega')$ are sampled, couple almost surely the trajectories touching $L$ using assumption $(S)$, so that they coincide (up to time shift) outside some random finite set $M$. Finally, take trajectories not touching $L$ to be the same in $\omega$ and in $\omega'$.
\end{remark}







\end{document}




****************************************************************************************************************************
Below is the bibliography of article with Serguei Popov.



\bibitem{A21}
\textsc{Y.~Abe} (2021) 
Second-order term of cover time for planar simple random 
walk.
\textit{J.\ Theoret.\ Probab.} \textbf{34} (3), 1689--1747.

% \bibitem{AZ22}
% \textsc{Y.~Abe, O.~Zeitouni} (2022)
% Private communications.

\bibitem{BK14}
\textsc{D.~Belius, N.~Kistler} (2017)
The subleading order of two dimensional cover times.
\textit{Probab.\ Theory Relat.\ Fields} \textbf{167} (1), 
461--552

\bibitem{BGP18}
\textsc{D.F.~de Bernardini, C.F.~Gallesco, S.~Popov} (2018)
On uniform closeness of local times of Markov chains 
and i.i.d.\ sequences.
\textit{Stochastic Process.\ Appl.} \textbf{128} (10),
3221--3252.

% \bibitem{CGPV13} \textsc{F.~Comets, C.~Gallesco, S.~Popov,
% M.~Vachkovskaia} (2013)
% On large deviations for the cover time of two-dimensional 
% torus. 
% \textit{Electr.\ J.\ Probab.}, \textbf{18}, article 96. 

\bibitem{CP17} \textsc{F.~Comets, S.~Popov} (2017)
The vacant set of two-dimensional critical
 random interlacement is infinite.
\textit{Ann.\ Probab.} \textbf{45} (6B), 4752--4785.
 
\bibitem{CP20} \textsc{F.~Comets, S.~Popov} (2020) 
  Two-dimensional Brownian random interlacements.
\textit{Potential Analysis}, \textbf{53} (2), 727--771. 
 
\bibitem{CPV16} \textsc{F.~Comets, S.~Popov, M.~Vachkovskaia} (2016)
Two-dimensional random interlacements and late points 
for random walks.
\textit{Commun.\ Math.\ Phys.} \textbf{343}, 129--164.

\bibitem{LL10} 
\textsc{G.~Lawler, V.~Limic} (2010)
\textit{Random Walk: A Modern Introduction.} 
Cambridge Studies in Advanced Mathematics, \textbf{123}. 
Cambridge University Press, Cambridge.
 
\bibitem{P21}
\textsc{S.~Popov} (2021)
\textit{Two-dimensional Random Walk: 
From Path Counting to Random Interlacements.}
Cambridge University Press, Cambridge. 

\bibitem{SLT} 
\textsc{S.~Popov, A.~Teixeira} (2015)
Soft local times and decoupling of random interlacements. 
\textit{J.\ European Math. Soc.} \textbf{17} (10), 2545--2593.

\bibitem{ProTyk11}
\textsc{E.~Procaccia, J.~Tykesson} (2011)
Geometry of the random interlacement.
\textit{Elect.\ Commun.\ Probab.} \textbf{16}, 528--544.

\bibitem{Rod19}
\textsc{P.-F.~Rodriguez} (2019)
 On pinned fields, interlacements, and random walk 
 on $(\mathbb{Z}/N\mathbb{Z})^2$.
\textit{Probab.\ Theory Relat.\ Fields}, \textbf{173} (3),
1265--1299.

\bibitem{S10} \textsc{A.-S.~Sznitman} (2010)
Vacant set of random interlacements and percolation.
\textit{Ann.\ Math.\ (2)}, \textbf{171} (3), 2039--2087.


